% Korean Noether paper 04 — mechanically assembled inherited production body.
% Origin reopened against ED0004; current scope-rebase authority: ED0011 / NOETH-DE-AUTH-v047-20260808.

% inherited: noether_paper04_ko_translation_001_20260804\targets\Noether_P04_Korean_T01_U01_UNCHECKED.tex
% sha256: 64297D37AC0E60F2E622A3231BA5B59846A1D00CB40A95D865B864DC9A30D2F1
\begin{center}
{\Large\bfseries $n$개 변수의 형식에 관한 불변식론\footnote{잘츠부르크 자연과학자 대회에서 발표한 짧은 보고를 참조하라. \emph{Jahresbericht d. deutsch. Math.-Vereinigung}, Bd.\ 19, S.\ 101.}}\par
\vspace{0.5em}
에를랑겐의 \emph{Emmy Noether} 양.\par
\medskip
\rule{0.22\linewidth}{0.4pt}\par
\vspace{0.7em}
Journal f.\ d.\ reine u.\ angew.\ Math.\ 139 (1911), S.\ 118--154
\end{center}

\begin{center}
\textbf{서론.}
\end{center}
$n$개 변수의 형식에 관한 사영 불변식론에서는, 이항 및 삼항 영역을 넘어서면 주된 문제, 즉 형식계의 유한성 증명에 뒤따르는 문제들이 거의 다루어지지 않았다. 그것은 형식들 사이의 상호 의존성과 이 연관관계를 다루는 방법에 관한 문제들이다. 이 방법들은 본질적으로 --- 이항 및 삼항 영역에서는 이것이 완전히 증명되어 있듯이 --- Study가 \glqq{}기호적 방법의 기본정리\grqq{}라고 부른 두 정리, 곧 형식의 기호적 표현 가능성에 관한 정리와 기호 항등식에 관한 정리에 기초한다. 뒤의 정리는 정수 불변 구성식들 사이에 성립하는 모든 관계를 소수의 기호 항등식을 차례로 적용하여 얻는다는 것, 따라서 기호적 방법만이 불변식의 영역을 벗어나지 않고 형식들의 연관관계를 파악하게 한다는 것을 말한다.\footnote{E. Study, \emph{Methoden zur Theorie der ternären Formen} (Leipzig, Teubner, 1889). II. \S\ 6. --- 특수 변환군의 불변식론에서도 이 기본정리들이 나타나는 것에 관해서는 Study, \emph{Das Apollonische Problem} (Math. Ann. Bd. 49 [1897]) 및 \emph{Zur Differentialgeometrie der analytischen Kurven} (Transactions of the American Math. Society, Bd. 1 [1909])을 참조하라.}

$n$개 변수로 넘어갈 때 생기는 어려움의 근거는 Gordan이 이미 자신의 에를랑겐 강령에서 지적하였다.\footnote{P. Gordan, \emph{Über das Formensystem binärer Formen}. S. 46. (Leipzig, Teubner, 1875).} 그것은 기호적 방법의 첫째 정리가 그 일반적 형태로는 더 이상 성립하지 않는 데 있다. 잘 알려져 있듯이 $n$항 영역에서는 더 높은 단계의 변수열 $p_\rho$\footnote{표기는 \S\ 1을 참조하라.}가 나타나며, 이와 마찬가지로 더 높은 단계의 기호열 $S^\rho$도 나타난다. 실제로 처음 Clebsch가 했듯이\footnote{A. Clebsch, Über eine Fundamentalaufgabe der Invariantentheorie. (Abh. der Ges. d. Wiss. zu Göttingen, Bd. XVII, 1872).} 임의의 개수의 열 $p_\rho$를 지닌 형식에서 출발하든, F. Deruyts와 A. Capelli의 방식에 따라 서로 동변인 임의의 개수의 열 $x$만을 지닌 형식에서 출발하든, 이 열들에 도달한다.\footnote{A. Capelli, \emph{Lezioni sulla teoria delle forme algebriche} (Napoli 1902), \S\ 22--24 및 거기에 제시된 문헌을 참조하라.} 그런데 열 $p_\rho,S^\rho$를 명시적으로\footnote{명시적 표현의 더 정확한 정의는 \S\ 2에 있다.} 포함하는 기호적 표현은 존재하지 않는다. 따라서 $p_\rho,S^\rho$를 각각의 열 $x$와 이에 반변인 열 $s$로 분해하고, 이들을 여러 행렬식에 분배해야 한다. 물론 미분 조건(공식 (19.) 참조)을 이용하면 주어진 행렬식 결합식이 오직 열 $p_\rho,S^\rho$에만 의존한다는 성질을 검증할 수 있다. 그러나 이 방법으로는 불변 구성식 전체와 그것을 생성하는 과정 전체를 조망할 수 없다.\footnote{Study에서 유래한, 계산상 매우 유용한 기호법의 한 발전(전달: F. Engel, Ber. d. K. Sächs. Ges. d. Wiss., math.-phys. Kl., 1900, S. 126; 4항 영역에 대해서는 E. Study, \emph{Geometrie der Dynamen}, S. 135)도 여기서 말하는 의미의 명시적 표현은 아니다. 예컨대 그것으로부터 형식을 생성하는 비기호적 미분 과정을 얻기는 어려울 것이다.}

이제 우리는 \S\S\ 2와 6에서 \glqq{}대응 행렬에 관한 정리\grqq{}\footnote{1862년 Graßmann의 확장론 Nr. 112를 제외하면, A. Brill, Über zwei Berührungsprobleme, \S\ 2 (Math. Ann. Bd. 4. 1871)에 처음 나타난다.}를 적용하여, 행렬식에 의한 표현을 \glqq{}행렬곱\grqq{}에 의한 표현으로 바꿈으로써 첫째 기본정리를 일반적인 형태로 다시 얻는 방법을 보인다. \S\ 2에서는 열 $S^\rho,p_\rho$를 명시적으로 포함하는 모든 행렬곱이 기본형의 불변 구성식임을 보인다. \S\ 6에 주어진 역, 즉 모든 불변 구성식이 행렬곱으로 명시적으로 표현된다는 사실은 둘째 기본정리를 전이한 뒤에야(\S\S\ 3--5) 얻어진다.

이 정리는 변수열 $x$만을 포함하는 형식에 대해서는 알려져 있다.\footnote{E. Pascal, Giorn. di mat. 26 (1888), S. 33, 102; Rendic. d. Accad. d. Linc. (4) 4 (1888), S. 119; ib. Mem. (4) 5 (1888), S. 375.} 모든 열 $S^\rho,p_\rho$를 분해 불가능한 것으로 보았을 때의 분해 불가능한 항등식 전체를 세우는 일반 문제는 Study가 정식화하였다.\footnote{\glqq{}Methoden \dots\grqq{}, S. 204, Anmerkung 17).} 그는 동시에 나타나는 항등식의 개수를 $n$항 영역에서 방법의 어려움을 재는 척도로 본다. 이제 항등식 전체를 하나의 공식 (36.)에 모을 수 있으므로, 이 어려움은 적어도 최소한으로 줄어든다. 다만 실제 적용에서는 (36.)의 몇몇 두드러진 특수한 경우가 자주 나타난다. 예컨대 새 열을 대입하지 않고도 명시적 표현을 허용한다는 특징을 지닌 (20.), (21.), 또는 열 $p_\rho$가 겉보기에는 열 $y$들로 분해되어 나타나므로 \glqq{}분해 항등식\grqq{}이라 부르는 공식 (31.)이 그러하다.

이 두 기본정리 위에서 나머지 이론은 쉽게 세울 수 있다. 첫째 정리는 형식을 생성하는 과정, 즉 기호적 수축 과정과 비기호적 미분 과정을 직접 준다.\footnote{\glqq{}Invariante Gebilde von Nullsystemen\grqq{} (Sitzungsber. d. Ak. d. Wiss. Wien, Bd. 97, Abt. IIa, 1888)에서 Mertens는 다른 방식, 곧 직접 일반화할 수 없는 방식으로 4항 수축 과정에 도달하였다. 거기서 나타나는 여섯 열 $S^2$의 교대 불변식은 여기서 대입 (11.)을 한 번 적용하여 생기는 불변식과 짝수 불변식의 한 결합식만큼 다르다. 특수한 한 경우의 4항 수축 과정은 P. Gordan, \emph{Das simultane System von zwei quadratischen quaternären Formen} (Math. Ann. Bd. 56. 1903)에도 나타난다.} 한편 분해 항등식은 이 과정들 가운데 서로 동치인 것을 모두 제거한다(\S\ 6). 더 나아가 미분 과정을 알면 \glqq{}정규형\grqq{}이라는 개념을 $n$항 영역으로 옮길 수 있고, Mertens가 4항 영역에서 제시한 증명법\footnote{F. Mertens, Über invariante Gebilde quaternärer Formen. (Sitzber. d. Ak. d. Wiss. Wien. Bd. 98. Abt. IIa. 1889.)}을 써서 불변 구성식의 한 계를 그와 동치인 정규형의 계로 바꿀 수 있다는 정리를 얻는다(\S\ 7). 나는 이러한 마지막 가정 아래에서, 삼항 수축 과정을 두 기본 수축을 차례로 적용하여 생성할 수 있음을 학위논문에서 보였다.\footnote{Über die Bildung des Formensystems der ternären biquadratischen Form. (Dieses Journal, Bd. 134. 1908.)} 이 정리와 급수 전개 이론에 기초하여 \glqq{}형식열\grqq{} 개념을 도입하고 형식계의 주요 환원정리를 전개하였다. 완전히 마찬가지로 $n$항 영역에서는 수축 과정을 $n-1$개의 기본 수축으로부터 생성할 수 있고(\S\ 8), 환원정리를 세울 수 있다(\S\ 9).

\paragraph{후기.} 잘츠부르크 보고 때 빈의 E. Müller 교수는 이 논문의 바탕을 이루는 행렬곱 계산이 원리상 Graßmann의 확장론 계산법과 동일하다는 점을 내게 알려 주었다.\footnote{Hermann Graßmanns gesammelte mathematische und physikalische Werke. Ersten Bandes zweiter Teil: Die Ausdehnungslehre von 1862. In Gemeinschaft mit H. Graßmann d. Jüngeren herausgegeben v. Fr. Engel (Leipzig, Teubner 1896).} 실제로 Graßmann의 \glqq{}조합곱\grqq{} $[S^{\rho_1}S^{\rho_2}\cdots S^{\rho_i}]$은 이 열들로 주어진 행렬로 볼 수 있다(\S\ 2 참조). 이때 \glqq{}전진곱\grqq{}은 열 자체의 행렬에, \glqq{}후퇴곱\grqq{}은 대응하는 열 $S_{n-\rho_1},S_{n-\rho_2},\ldots,S_{n-\rho_i}$의 행렬에 해당한다. \glqq{}혼합곱\grqq{}에 해당하는 행렬은 여러 열 $S$를 대입 (9.)로 새 열 $P,Q$에 모아 만들게 된다. 여기서 우리의 \glqq{}대응 열\grqq{}은 Graßmann의 \glqq{}보충\grqq{}에, 우리의 \glqq{}행렬곱\grqq{}은 \glqq{}단계수가 0인 혼합곱\grqq{}에 해당한다.

이 대응 아래에서 1862년 확장론 Nr. 113에 주어진 전개는 우리의 공식 (32.)와 쌍대인 공식과 동일해진다. 특수한 경우 $\rho=1$에서 (32.)는 (17.)로, 쌍대 공식은 (16.)으로 바뀐다. 최근 Müller는 Graßmann 전개의 이 특수한 경우들로부터\footnote{E. Müller, Beiträge zur Graßmannschen Ausdehnungslehre. 1. Mitteilung: Einige allgemeine Sätze. (Sitzungsber. d. Ak. d. Wiss. Wien; Bd. 118. Abt. IIa. Juli 1909), Satz 16 und 18.} 앞서 말한 항등식 (20.), (21.)을 유도하였다.

우리의 유도는 Graßmann--Müller의 유도와 달리 항등식의 완전성까지 동시에 증명한다. 여기서 다루는 문제에는 바로 이것이 가장 본질적이다. 또한 우리의 유도는 (20.), (21.)에서 출발하여 지금까지 주목되지 않은 듯한 가장 일반적인 분해 불가능 항등식 (36.)으로 나아간다.

% inherited: noether_paper04_ko_translation_001_20260804\targets\Noether_P04_Korean_T01_U02_UNCHECKED.tex
% sha256: 2728D4B1AA51C0CE3BCAC45D274BFD9B9ACB2F96DA4AFF2518F1CD202A287CC5
\section*{\S\ 1. 표기와 정의. 알려진 결과의 요약.}
통상적인 방식대로 원소 $x_1,x_2,\ldots,x_n$의 열을 변수열 $x$라 부른다. 이와 마찬가지로 변수열 $p_\rho$ ($\rho=1,2,\ldots,n-1$)라 함은 $\binom n\rho$개의 원소 $p_{i_1i_2\ldots i_\rho}$로 이루어진 열을 뜻한다. 여기서 $p_{i_1i_2\ldots i_\rho}$는 $\rho$개의 열 $x^{(1)},x^{(2)},\ldots,x^{(\rho)}$의 행렬에서 얻는 다음 $\binom n\rho$개의 행렬식이다.
\[
\left|\begin{array}{cccc}
 x_{i_1}^{(1)}&x_{i_2}^{(1)}&\cdots&x_{i_\rho}^{(1)}\\
 x_{i_1}^{(2)}&x_{i_2}^{(2)}&\cdots&x_{i_\rho}^{(2)}\\
 \vdots&\vdots&&\vdots\\
 x_{i_1}^{(\rho)}&x_{i_2}^{(\rho)}&\cdots&x_{i_\rho}^{(\rho)}
\end{array}\right|,
\qquad (i_1<i_2<\cdots<i_\rho=1,2,\ldots,n).
\]
대응 행렬에 관한 정리\footnote{예를 들어 Gordan--Kerschensteiner, \emph{Vorlesungen über Invariantentheorie}, Bd. I, S. 99를 참조하라.}에 따르면, 각 열 $p_\rho$에는 그 열의 각 원소마다 성립하는 다음 관계를 통하여 둘째 열 $q_{n-\rho}$가 일대일로 대응한다.
\begin{equation}
 p_{i_1i_2\ldots i_\rho}=(-1)^{\frac{\rho(\rho+1)}{2}+i_1+i_2+\cdots+i_\rho}q_{j_1j_2\ldots j_{n-\rho}},
\tag{1}
\end{equation}
여기서 $i$들과 $j$들은 각각 오름차순으로 놓여 있고, 이들을 합하면 1부터 $n$까지의 수를 한 번씩 모두 갖는 순열이 된다. 열 $q_{n-\rho}$는 $(n-\rho)$차 행렬식 $q_{j_1j_2\ldots j_{n-\rho}}$ $\binom n\rho$개로 이루어진다. 이 행렬식들은 $x$에 반변인 $n-\rho$개의 열 $u^{(1)},u^{(2)},\ldots,u^{(n-\rho)}$의 행렬에서 만들어진다. 이 열들은 다음 $\rho(n-\rho)$개의 조건방정식을 만족한다.
\begin{equation}
 (u^{(k)}\mid x^{(l)})=\sum_{\alpha=1}^{n}u_\alpha^{(k)}x_\alpha^{(l)}=0;
 \qquad (k=1,2,\ldots,n-\rho;\ l=1,2,\ldots,\rho).
\tag{2}
\end{equation}
우리는 Clebsch를 따라\footnote{a. a. O. \S\ 5.} 이 열들을 정규화하여, (1.)에 원래 들어가는 비례인자를 1로 놓기로 한다.

% inherited: noether_paper04_ko_translation_001_20260804\targets\Noether_P04_Korean_T01_U03_UNCHECKED.tex
% sha256: 778BA22BFDE7F0C7A66F1CC6B0FB1A64B3B864F81BEE4AB78CF6B9735B14A733
Clebsch에 따르면\footnote{a. a. O. \S\ 12.}, 그리고 뒤에 나온 F. Deruyts와 A. Capelli의 연구에 따르면(서론 참조), 가장 일반적인 형식들은 $n-1$개의 각 열 $p_\rho$를 기껏해야 하나씩만 포함하는 형식들로 환원할 수 있다. 따라서 이들은 기호적으로 다음과 같이 표현된다.
\begin{equation}
 (S^1\mid p_1)^{m_1}(S^2\mid p_2)^{m_2}\cdots(S^{n-1}\mid p_{n-1})^{m_{n-1}},
\tag{3}
\end{equation}
여기서 (2.)와 유사하게 다음과 같이 놓았다.
\[
 (S^\rho\mid p_\rho)=\sum S^{i_1i_2\ldots i_\rho}p_{i_1i_2\ldots i_\rho}.
\]
한 열 $p_\rho$의 원소들 사이에 성립하는 비선형 항등관계 때문에, 기호열 $S^\rho$는 바로 그 항등관계들을 만족하도록 정규화할 수 있다. 즉 이 열은 $x$에 반변인 열 $s^{(1)},s^{(2)},\ldots,s^{(\rho)}$로 이루어진 $\rho$행 행렬의 행렬식처럼 행동한다.\footnote{Clebsch, a. a. O. \S\ 4.} 따라서 각 열 $S^\rho$에는 다음 관계로 다른 한 열 $S_{n-\rho}$가 일대일로 대응한다.
\begin{equation}
 S_{i_1i_2\ldots i_{n-\rho}}=(-1)^{\frac{(n-\rho)(n-\rho+1)}2+i_1+i_2+\cdots+i_{n-\rho}}S^{j_1j_2\ldots j_\rho},
\tag{4}
\end{equation}
여기서 다시 $i$들과 $j$들은 각각 오름차순으로 놓여 있으며, 합쳐서 1부터 $n$까지의 수를 정확히 모두 소진한다. 열 $S_{n-\rho}$의 원소 $S_{i_1i_2\ldots i_{n-\rho}}$는 $x$에 동변인 $n-\rho$개의 열 $\sigma^{(1)},\sigma^{(2)},\ldots,\sigma^{(n-\rho)}$로 만든 행렬식처럼 행동하며, 열 $s$와 $\sigma$는 다음 $\rho(n-\rho)$개의 조건으로 연결된다.
\begin{equation}
 (s^{(k)}\mid \sigma^{(l)})=0,\qquad(k=1,2,\ldots,\rho;\ l=1,2,\ldots,n-\rho).
\tag{5}
\end{equation}
관계 (1.)과 (4.)에 따라 (3.)의 각 인자는 다음 네 가지 서로 동치인 기호적 표현을 갖는다.
\begin{equation}
 (S^\rho\mid p_\rho)=(S^\rho q_{n-\rho})=(S_{n-\rho}p_\rho)=(-1)^{\rho(n-\rho)}(q_{n-\rho}\mid S_{n-\rho}\footnotemark),
\tag{6}
\end{equation}
\footnotetext{Clebsch, a. a. O. \S\ 5를 참조하라.}
이하에서 언제나 그렇듯이, $(S^\rho q_{n-\rho})$와 $(S_{n-\rho}p_\rho)$는 각각 행렬식
$(s^{(1)}\cdots s^{(\rho)}u^{(1)}\cdots u^{(n-\rho)})$와
$(\sigma^{(1)}\cdots\sigma^{(n-\rho)}x^{(1)}\cdots x^{(\rho)})$의 약기호이다. 라플라스 전개를 하면 이 행렬식들이 각각 열 $S^\rho,q_{n-\rho}$ 및 $S_{n-\rho},p_\rho$에만 의존한다는 것이 드러나며, 위 표기는 바로 이것을 나타내려는 것이다. 앞의 설명에 따라 $(S^\rho\mid p_\rho)$와 $(q_{n-\rho}\mid S_{n-\rho})$는 각각 열 $s$와 $x$의 행렬, 그리고 열 $u$와 $\sigma$의 행렬을 곱한 것을 뜻한다. 여기서 행렬곱이란 서로 대응하는 행렬식들의 곱을 모두 더한 것, 즉 곱행렬의 행렬식 값을 말한다.

% inherited: noether_paper04_ko_translation_001_20260804\targets\Noether_P04_Korean_T01_U04_UNCHECKED.tex
% sha256: 578FFBC95675FDD5AA822287833EA3CBBDB53752FC8FEAAE85F7B156578F1F59
기호열(또는 변수열)에 속하는 특성수 $\rho(n-\rho)$를 Clebsch를 따라 그 열의 무게라 부르고,\footnote{a. a. O. \S\ 11.} 열 $s$의 개수를 나타내는 수 $\rho$를 위 무게지수, 열 $\sigma$의 개수를 나타내는 수 $n-\rho$를 아래 무게지수라 부른다. 관계 (4.)에 따르면 하나의 기호열은 같은 관계 안에서 한 번은 위 무게지수를, 한 번은 아래 무게지수를 달고 나타날 수 있다. 서로 짝을 이루는 열 $p_\rho$와 $q_{n-\rho}$에도 같은 사실이 성립한다. 또한 일반적으로 다음과 같이 표기하기로 한다. 대응 행렬이 열 $x$들로 이루어진 더 높은 단계의 변수열은 $p$로, 그 행렬이 열 $u$들로 이루어진 것은 $q$로 나타낸다. $x$에 동변인 기호열은 작은 그리스 문자 $\sigma,\tau,\alpha,\beta$로, $u$에 동변인 기호열은 작은 라틴 문자 $s,t,a,b$로 나타낸다. 그 밖의 모든 기호열은 무게지수를 붙인 큰 라틴 문자 $S^\rho,T^\tau,A^\rho$ 또는 $S_{n-\sigma},T_{n-\tau},A_{n-\rho}$로 나타낸다. 위 무게지수 또는 아래 무게지수에 맞추어, (4.)에서처럼 한 열의 개별 원소도 위 첨자 또는 아래 첨자로 쓴다.

% inherited: noether_paper04_ko_translation_001_20260804\targets\Noether_P04_Korean_T02_U05_UNCHECKED.tex
% sha256: 32224493EE78CF4B206732EBCAD5BABCB6D9F2981D52BAFDB6C09878ADE53C06
\section*{\S\ 2. 행렬곱에 의한 표현.}
이하에서 우리는 언제나 ``행렬곱''을, 같은 수의 열로 이루어진 두 행렬에서 서로 대응하는 행렬식들의 곱을 모두 합한 것으로 이해한다. 이때 둘째 행렬의 열들은 첫째 행렬의 대응 열들에 반변한다. 각 행렬의 열들은 1. 하나씩 주어질 수도 있고, 2. (6.)에서처럼 하나의 열 $S^\rho$로 묶일 수도 있으며, 3. 서로 다른 여러 열 $S^\rho,S^\sigma\ldots$로 합쳐질 수도 있다. 행렬곱은 (6.)에서 사용한 기호 $(\ \mid\ )$로 나타내기로 한다. 표현 (6.)은 대응 행렬에 관한 정리의 둘째 정식화를 특수하게 적용한 것이다. ``모든 $\rho$열 행렬곱 $(v^{(1)}v^{(2)}\cdots v^{(\rho)}\mid y^{(1)}\cdots y^{(\rho)})$에는 하나의 행렬식이 대응한다. 역으로 모든 행렬식에는 미리 정한 열의 수 $\rho$ ($\rho=1,2,\ldots,n-1$)를 가진 하나의 행렬곱을 대응시킬 수 있다.''\footnote{$\rho>n$이면 행렬곱은 잘 알려져 있듯이 소멸하고, $\rho=n$이면 두 행렬식의 곱이 된다.} 실제로 주어진 행렬곱에서 행렬식으로, 그리고 그 역으로 넘어가기 위해서는 열 $y$들을 하나의 열 $p_\rho$로 묶거나(각각 열 $v$들을 하나의 열 $q_\rho$로 묶고) (1.)에 주어진 대입을 실행하기만 하면 된다. 다만 이때 행렬식의 값은 그 개별 원소들에 의해 직접 주어지는 것이 아니라 라플라스 전개를 거쳐서야 얻어진다. 이와 같이 행렬식과 행렬곱이 동치이므로, 형식의 기호적 표현 가능성에 관한 정리(첫째 기본정리)는 다음과 같이도 표현된다.

\emph{정리 I:} ``모든 불변 구성식은 행렬곱으로 기호적으로 표현할 수 있으며, 역으로 모든 행렬곱은 불변 구성식이다.''\footnote{미리 정한 기본형들의 불변 구성식은 물론 기호열 $S^\rho\ldots$에 의존하는 행렬곱의 결합식뿐이다.}

행렬곱에 의한 표현을 쓰면, 행렬식 표현에 따르는 본질적 난점, 곧 표현이 명시적이지 않다는 난점을 극복할 수 있다.\footnote{서론을 참조하라.} 여기서 ``명시적 표현''이란 열 $S^\rho,T^\tau,\ldots$ 자체나 이들로부터 유일하게 유도할 수 있는 열 $R^\rho$만이 들어가고, 개별 부분열 $s,t,\ldots$는 더 이상 나타나지 않는 표현을 뜻한다. 우리는 \S\ 6에서 기호열 $S^\rho,T^\tau,\ldots$에만 의존하는 모든 불변 구성식에 대하여 이 기호열들만으로 된 명시적 표현을 얻고, 그와 함께 생성 과정인 수축 과정과 미분 과정을 $n$항 영역으로 옮길 것이다. 명시적 표현이 가능하다는 사실은 다음에서 분명해진다. 가장 단순한 행렬곱들만이 행렬식 표현에서 나타나는 행렬식들에 그대로 대응할 수 있으므로, 그 밖의 모든 행렬곱은 서로 다른 여러 행렬식을 하나의 식으로 묶은 것이다. 실제로 기호열과 변수열을 개별 열 $s,u,x$로 풀어 놓으면, 첫째 기본정리의 통상적인 형태에 따라 반드시 행렬식의 결합식이 생긴다. 역으로, 미리 정한 기본형들의 불변 구성식을 나타내는 모든 행렬식 결합식을 명시적 행렬곱으로 묶을 수 있다는 증명은 둘째 기본정리를 사용하여 \S\ 6에서 얻는다.

% inherited: noether_paper04_ko_translation_001_20260804\targets\Noether_P04_Korean_T02_U06_UNCHECKED.tex
% sha256: 8EFCEE86349CBB92A6BAF9A5A607E516F2EADF26FC5A82462E0F8FE78BBD574A
이제 두 기호열 $S^\sigma,T^\tau$와 변수열을 이용하여 기본형 $(S^\sigma\mid p_\sigma)(T^\tau\mid p_\tau)$에서 모든 열을 명시적으로 포함하는 불변 구성식을 얻으려면(수축 과정), 이 절의 첫머리에 제시한 정의에 따라 제3종 행렬곱을 하나 만들기만 하면 된다.
\begin{equation}
 \pair{q_{n-\sigma-\lambda}S^\sigma}{T_{n-\tau}p_{\tau-\lambda}},
 \qquad (\lambda=1,2,\ldots,\tau,n-\sigma\footnotemark).
\tag{7}
\end{equation}
\footnotetext{수 $\lambda$는 언제나 마지막에 적은 두 수 가운데 작은 것까지만 변한다.}
(7.)을 전개하면 다음을 얻는다.
\begingroup
\setlength{\arraycolsep}{2pt}
\begin{equation}
\begin{aligned}
&\pair{v^{(1)}\cdots v^{(n-\sigma-\lambda)}s^{(1)}\cdots s^{(\sigma)}}
       {\alpha^{(1)}\cdots\alpha^{(n-\tau)}y^{(1)}\cdots y^{(\tau-\lambda)}} \\
&=\sum
\left|\begin{array}{ccc}
 v_{i_1}^{(1)}&\cdots&v_{i_{n-\lambda}}^{(1)}\\
 \vdots&&\vdots\\
 v_{i_1}^{(n-\sigma-\lambda)}&\cdots&v_{i_{n-\lambda}}^{(n-\sigma-\lambda)}\\
 s_{i_1}^{(1)}&\cdots&s_{i_{n-\lambda}}^{(1)}\\
 \vdots&&\vdots\\
 s_{i_1}^{(\sigma)}&\cdots&s_{i_{n-\lambda}}^{(\sigma)}
\end{array}\right|
\left|\begin{array}{ccc}
 \alpha_{i_1}^{(1)}&\cdots&\alpha_{i_{n-\lambda}}^{(1)}\\
 \vdots&&\vdots\\
 \alpha_{i_1}^{(n-\tau)}&\cdots&\alpha_{i_{n-\lambda}}^{(n-\tau)}\\
 y_{i_1}^{(1)}&\cdots&y_{i_{n-\lambda}}^{(1)}\\
 \vdots&&\vdots\\
 y_{i_1}^{(\tau-\lambda)}&\cdots&y_{i_{n-\lambda}}^{(\tau-\lambda)}
\end{array}\right|\\
&=\sum_{i,k,l}\eps\,q_{k_1\ldots k_{n-\sigma-\lambda}}
 S^{k_{n-\sigma-\lambda+1}\ldots k_{n-\lambda}}
 T_{l_1\ldots l_{n-\tau}}
 p^{l_{n-\tau+1}\ldots l_{n-\lambda}},
\end{aligned}
\tag{8}
\end{equation}
\endgroup
즉, 이는 열 $S,T,q,p$에만 의존하는 식이다. 여기서 $\eps=\pm1$이고\footnote{이 표기는 이하 줄곧 유지한다.}, 합은 다음과 같이 이해한다. 각 열 $S\ldots$의 지수들은 그 열 안에서 올바른 순서로 배열하고, 각 조합 $i_1,i_2,\ldots,i_{n-\lambda}$에 대하여 $k$와 $l$은 각각 그때 가능한 수 $i$들의 모든 순열을 따로 돈다. $\sigma\geqq\tau$인 경우 (7.)에 들어가는 수 $\lambda$를 수축의 결손이라 부른다. 이 마지막 제한의 이유는 \S\ 6에서 드러난다.

% inherited: noether_paper04_ko_translation_001_20260804\targets\Noether_P04_Korean_T02_U07_UNCHECKED.tex
% sha256: 8F03D5DA8AC825B5F180049C1AC9CB01DB52EE7507F2FAE60471744ADF544A7B
(7.)에 들어가는 두 행렬의 행렬식들은 새로운 열 $Q^{n-\lambda},P_{n-\lambda}$의 원소로 볼 수 있으며, 이들에 대응하는 열을 $Q_\lambda,P^\lambda$라 하자. 이 열 $Q,P$는 (8.)에 따라 원래의 $S$와 $q$, 또는 각각 $T$와 $p$로 나타낼 수 있다. 이를 다음 등식으로 표시한다.
\begin{equation}
\begin{aligned}
 q_{n-\sigma-\lambda}S^\sigma&=Q^{n-\lambda}\sim Q_\lambda, &\quad\text{또는}\quad Q_\lambda&=[q_{n-\sigma-\lambda}S^\sigma]_\lambda,\\
 T_{n-\tau}p_{\tau-\lambda}&=P_{n-\lambda}\sim P^\lambda, &\quad\text{또는}\quad P^\lambda&=[T_{n-\tau}p_{\tau-\lambda}]^\lambda.
\end{aligned}
\tag{9}
\end{equation}
(4.)를 고려하면, (6.)과 마찬가지로 행렬곱 (7.)에는 다음과 같은 서로 동치인 네 가지 기호적 표현이 성립한다.
\begin{equation}
\pair{q_{n-\sigma-\lambda}S^\sigma}{T_{n-\tau}p_{\tau-\lambda}}
=(q_{n-\sigma-\lambda}S^\sigma P^\lambda)=(Q_\lambda T_{n-\tau}p_{\tau-\lambda})=(-1)^{\lambda(n-\lambda)}\pair{P^\lambda}{Q_\lambda}.
\tag{10}
\end{equation}
(7.)에는 다음과 같이 놓음으로써 새로운 열을 한 번 더 대입할 수 있다.
\begin{equation}
 \pair{q_{n-\sigma-\lambda}S^\sigma}{T_{n-\tau}p_{\tau-\lambda}}
 =\pair{R^{\sigma+\lambda}}{p_{\sigma+\lambda}}(R^{\tau-\lambda}p_{\tau-\lambda}).
\tag{11}
\end{equation}
여기서도 (8.)은 열 $R$의 원소들의 곱을 열 $S$와 $T$로 유일하게 나타낼 수 있음을 보여 준다. 대입 (9.)와 (11.)은 특히 (7.)에 다른 열들을 사용한 수축을 한 번 더 시행할 때 적용하게 된다.

% inherited: noether_paper04_ko_translation_001_20260804\targets\Noether_P04_Korean_T02_U08_UNCHECKED.tex
% sha256: 89DDC2F775CD86A216346C080B45E81527B60C1A8F115F107842D477F1E8D857
행렬 표현을 이용하면 열 $p_\rho$의 원소들 사이의 이차 관계(잘 알려져 있듯이 나머지 모든 관계가 여기서 따른다), 따라서 \S\ 1에 따르면 기본형의 계수에 관해 선형인, 열 $S^\rho$가 만족하는 관계를 불변적으로 세울 수 있다. 그것은 다음 항등식들이다.
\begin{equation}
 \pair{q_{n-\rho-\lambda}S^\rho}{S_{n-\rho}p_{\rho-\lambda}}=0\footnotemark,
 \qquad (\lambda=1,2,\ldots,\rho,n-\rho),
\tag{12}
\end{equation}
\footnotetext{Graßmann의 저작 제1권 제2부, S. 510에 있는 Study의 후기를 참조하라($S^\rho$라는 양이 단순하기 위한 조건).}
여기서 $p$와 $q$는 임의의 양으로 보아야 한다. 결손수 $\lambda$가 홀수인 항등식들은 더 높은 결손의 항등식들로부터 따른다는 것을 \S\ 5에서 볼 것이다. 관계 (12.)는 관계 (5.)의 직접적인 귀결이다. 실제로
\[
 \pair{q_{n-\rho-\lambda}S^\rho}{S_{n-\rho}p_{\rho-\lambda}}
 =(v^{(1)}\cdots v^{(n-\rho-\lambda)}s^{(1)}\cdots s^{(\rho)}\mid \sigma^{(1)}\cdots\sigma^{(n-\rho)}y^{(1)}\cdots y^{(\rho-\lambda)}),
\]
이고, 곱정리를 적용하면 (5.) 때문에 소멸하는 $(n-\lambda)$열 행렬식을 얻는다. 항등식 (12.)이 말하는 것은, 모든 수축 유형을 찾기 위해서는 각 변수열에 관해 선형인 형식들의 곱을 수축하는 것으로 충분하다는 것이다. 또는 \S\ 6에 따르면, 정규화된 형식 $(S^\rho\mid p_\rho)^m$은 계수에 관해 선형인 불변 구성식을 갖지 않는다는 것이다.

% inherited: noether_paper04_ko_translation_001_20260804\targets\Noether_P04_Korean_T02_U09_UNCHECKED.tex
% sha256: 5A940A5CECFD77AE33FCA9C2E3A713653D682886E9F5CD0FD66605DCD75B82E3
조건 (12.)는 열 $p_\rho$를 특징짓는 데에도 충분하다. 실제로 $p_\rho$의 개별 원소가 어떤 행렬의 행렬식들이라는 성질은 불변적이므로, 그 성질도 불변적으로 표현할 수 있어야 한다. 그런데 정리 VI에 따르면 조건 (12.)는 $p_\rho$에 가능한 가장 큰 불변적 제한을 나타낸다. 즉, 열 $p_\rho$를 계수로 삼아 만든 선형형식은 어떠한 불변 구성식도 갖지 않는다는 제한이다.

% inherited: noether_paper04_ko_translation_001_20260804\targets\Noether_P04_Korean_T03_U10_UNCHECKED.tex
% sha256: 66A5859EB5CA1C73D64E5116E720AF229CB26C70B0E060C3A94FB06A897E7E83
\section*{\S\ 3. 기호 항등식\footnote{문제의 정식화와 문헌에 대해서는 서론을 참조하라.}.}
이제 기호적 방법의 둘째 기본정리를 옮겨야 한다. 즉, 모든 열 $S^\rho,p_\rho$를 분해 불가능한 것으로 보았을 때의 기약인 분해 불가능 항등식 전체를 세워야 한다. 이를 위하여 다음을 추가로 정한다. 기호열과 변수열의 의미에 맞추어, 변수열 $p_\rho$를 $p_{\rho-1}y$로 특수화한 뒤 생기는 식에 이 계의 항등식 가운데 하나를 적용하여 얻는 항등식은 합성된 것으로 본다. 반면 $k$개의 기호열을 포함하는 항등식은, 위 과정으로 더 적은 기호열을 가진 항등식에서 생기더라도 분해 불가능한 것으로 본다. 열 $S^\rho,p_\rho$가 항등식에 반드시 명시적으로 들어갈 필요는 없다. 나타나는 행렬곱들이 이 열들에만 의존한다는 것을 보이면, 이들을 분해 불가능한 양으로 이해하기에 충분하다. 그러나 우리는 이 경우에도, 일부에는 대입 (11.)을 적용하여, 언제나 명시적 표현을 얻을 수 있음을 보일 것이다. 일반 항등식은 Pascal이 제시한 특수 항등식들로부터 행렬 표현의 원리를 써서 얻는다. 이 특수 항등식들은 열 $x$와 $u$만을 포함하며, Study가 삼항 영역에서 제시한 항등식을 직접 옮긴 것이다.
\begin{align}
\sum_{i=1}^n(-1)^{i+1}\pair{a^{(i)}}{x}(a^{(1)}\cdots a^{(i-1)}a^{(i+1)}\cdots a^{(n)}u)&=(-1)^{n+1}\pair{u}{x}(a^{(1)}\cdots a^{(n)}),\tag{13}\\
\sum\pm\pair{a^{(1)}}{x^{(1)}}\pair{a^{(2)}}{x^{(2)}}\cdots\pair{a^{(n)}}{x^{(n)}}&=(a^{(1)}\cdots a^{(n)})(x^{(1)}\cdots x^{(n)}),\tag{14}\\
\sum_{i=1}^n(-1)^{i+1}\pair{u}{a^{(i)}}(a^{(1)}\cdots a^{(i-1)}a^{(i+1)}\cdots a^{(n)}x)&=(-1)^{n+1}\pair{u}{x}(a^{(1)}\cdots a^{(n)}).\tag{15}
\end{align}
(13.)과 (15.)에 각각 $(a^{(i)}\mid x)=(a^{(i)}q_{n-1})$와 $(u\mid a^{(i)})=(p_{n-1}a^{(i)})$를 대입하면 두 항등식을 더 얻지만, 우리의 관점에서는 본질적으로 다른 것으로 보지 않는다. (14.)는 행렬식의 곱정리를 나타낸다. 이제 (13.)과 (14.)를(쌍대적으로는 (15.)와 (14.)를) 적절히 만든 행렬의 곱정리 계로 어떻게 바꿀 수 있는지 보이겠다. 곱정리를 한 번은 $\rho$열 행렬에, 한 번은 $(\rho-1)$열 행렬에 적용하면 다음을 얻는다.

% inherited: noether_paper04_ko_translation_001_20260804\targets\Noether_P04_Korean_T03_U11_UNCHECKED.tex
% sha256: 51533E6489F05B23B78B65A46C4AE131D6412680B6236D9F091421B4335CB67F
\[
\pair{a^{(1)}a^{(2)}\cdots a^{(\rho)}}{xp_{\rho-1}}
=(a^{(1)}a^{(2)}\cdots a^{(\rho)}\mid xy^{(2)}\cdots y^{(\rho)})
=\sum\pm\pair{a^{(1)}}{x}\pair{a^{(2)}}{y^{(2)}}\cdots\pair{a^{(\rho)}}{y^{(\rho)}},
\]
\[
\sum\pm\pair{a^{(2)}}{y^{(2)}}\cdots\pair{a^{(\rho)}}{y^{(\rho)}}
=(a^{(2)}\cdots a^{(\rho)}\mid y^{(2)}\cdots y^{(\rho)})
=(a^{(2)}\cdots a^{(\rho)}p_{\rho-1})=(a^{(2)}\cdots a^{(\rho)}q_{n-\rho+1});
\]
따라서 $\rho=2,3,\ldots,n$일 때 다음과 같은 곱정리 계가 따른다.
\begin{equation}
\begin{aligned}
\pair{a^{(1)}a^{(2)}\cdots a^{(\rho)}}{xp_{\rho-1}}
 &=\sum_{i=1}^{\rho}(-1)^{i+1}\pair{a^{(i)}}{x}(a^{(1)}\cdots a^{(i-1)}a^{(i+1)}\cdots a^{(\rho)}\mid p_{\rho-1})\\
 &=\sum_{i=1}^{\rho}(-1)^{i+1}\pair{a^{(i)}}{x}(a^{(1)}\cdots a^{(i-1)}a^{(i+1)}\cdots a^{(\rho)}q_{n-\rho+1}).\footnotemark
\end{aligned}
\tag{16}
\end{equation}
\footnotetext{항등식 (16.)은 $\rho$개 변수의 영역에서 가져온 항등식 (13.)으로도 볼 수 있다.}
$\rho=n$이면 (16.)은 (13.)이 된다. 더 나아가 $p_{n-1}=y^{(2)}p_{n-2}$, $p_{n-2}=y^{(3)}p_{n-3}$ 등으로 특수화하고, 식 $(a^{(1)}\cdots a^{(i-1)}a^{(i+1)}\cdots a^{(n)}\mid y^{(2)}p_{n-2})$, $(a^{(1)}\cdots a^{(i-1)}a^{(i+1)}\cdots a^{(h-1)}a^{(h+1)}\cdots a^{(n)}\mid y^{(3)}p_{n-3})$ 등에 항등식 (16.)을 적용하면 (13.)에서 (14.)에 도달한다. 따라서 우리의 정의에 따르면 항등식 (14.)는 항등식 (16.)들로부터 합성된 것이다. 달리 말해, 계 (16.)은 항등식 (13.) 및 (14.)와 동치이다. 계 (16.)은 일반 항등식에 요구한 조건 가운데 하나를 만족한다. 즉, 열 $p_{\rho-1}$을 분해 불가능한 양으로 포함하는 동시에 명시적인 형태로 포함한다. 이것은 이 요구만을 만족하는 유일한 계이다. 실제로 열 $a,x,p$로부터 다른 행렬식이나 행렬곱을 만들 수 없으며, 이들 사이에는 (16.)과 독립인 관계가 있을 수 없다. 그런 관계가 있다면 $(a^{(1)}\cdots a^{(\rho)}\mid xp_{\rho-1})$를 소거하여 $\rho$개($\rho\leqq n$)의 식 $(a^{(i)}\mid x)(a^{(1)}\cdots a^{(i-1)}a^{(i+1)}a^{(\rho)}q_{n-\rho+1})$ 사이의 관계를 얻을 것이다. 그러나 (13.)에 따르면 적어도 $n+1$개의 그러한 식이 하나의 관계에 들어가야 한다. 이와 유사한 고찰로부터 (15.) 및 (14.)와 동치인 다음 계를 얻는다.

% inherited: noether_paper04_ko_translation_001_20260804\targets\Noether_P04_Korean_T03_U12_UNCHECKED.tex
% sha256: AC6CD6F491B3EBAF83595B8BF7EDA17414679AF2E2D16B02CA1EFC59FC7589BA
\begin{equation}
\begin{aligned}
\pair{u q_{\rho-1}}{a^{(1)}\cdots a^{(\rho)}}
 &=\sum_{i=1}^{\rho}(-1)^{i+1}\pair{u}{a^{(i)}}(q_{\rho-1}\mid a^{(1)}\cdots a^{(i-1)}a^{(i+1)}\cdots a^{(\rho)})\\
 &=\sum_{i=1}^{\rho}(-1)^{i+1}\pair{u}{a^{(i)}}(p_{n-\rho+1}a^{(1)}\cdots a^{(i-1)}a^{(i+1)}\cdots a^{(\rho)}).
\end{aligned}
\tag{17}
\end{equation}

% inherited: noether_paper04_ko_translation_001_20260804\targets\Noether_P04_Korean_T03_U13_UNCHECKED.tex
% sha256: 953C46165395C344F8E1455994DE8956B13C427C790275B1C47E770BA5E6CE13
(16.)에서 임의의 열 $A^{\rho_1},B^{\rho_2},\ldots,x,p$ 사이의 항등식으로 넘어가기 위하여 다음을 생각한다. 이 항등식들은 $A^{\rho_1}=a^{(1)}a^{(2)}\cdots a^{(\rho_1)}$, $B^{\rho_2}=b^{(1)}b^{(2)}\cdots b^{(\rho_2)}$ 등으로 풀어 놓는 순간 (16.)과 같아져야 한다. 또한 개별 식들이 (16.)에 들어가는 유형이라면, (16.)으로 정해지는 연산들만을 통해서도 종결식으로 묶일 수 있어야 한다. 이제
\begin{equation}
 \rho=\rho_1+\rho_2+\cdots+\rho_k\leqq n
\tag{18}
\end{equation}
로 놓으면 다음을 얻는다.
\begin{align*}
&\pair{a^{(1)}\cdots a^{(\rho_1)}b^{(1)}\cdots b^{(\rho_2)}\cdots k^{(1)}\cdots k^{(\rho_k)}}{xp_{\rho-1}}
  =\pair{A^{\rho_1}B^{\rho_2}\cdots K^{\rho_k}}{xp_{\rho-1}} \\
&\quad=\sum_{i=1}^{\rho_1}(-1)^{i+1}\pair{a^{(i)}}{x}(a^{(1)}\cdots a^{(i-1)}a^{(i+1)}\cdots a^{(\rho_1)}B^{\rho_2}\cdots K^{\rho_k}q_{n-\rho+1})\\
&\qquad+(-1)^{\rho_1\rho_2}\sum_{i=1}^{\rho_2}(-1)^{i+1}\pair{b^{(i)}}{x}(b^{(1)}\cdots b^{(i-1)}b^{(i+1)}\cdots b^{(\rho_2)}A^{\rho_1}\cdots K^{\rho_k}q_{n-\rho+1})\\
&\qquad+\cdots\\
&\qquad+(-1)^{(\rho_1+\cdots+\rho_{k-1})\rho_k}\sum_{i=1}^{\rho_k}(-1)^{i+1}\pair{k^{(i)}}{x}(k^{(1)}\cdots k^{(i-1)}k^{(i+1)}\cdots k^{(\rho_k)}A^{\rho_1}\cdots (K-1)^{\rho_{k-1}}q_{n-\rho+1}).
\end{align*}
이제 각각의 합 전체는(그러나 그 합의 일부만으로는 안 된다) 실제로 열 $A^{\rho_1},B^{\rho_2}\ldots$를 분해 불가능한 양으로 포함한다. 그에 필요한 동시에 충분한 다음 조건들을 만족하기 때문이다.\footnote{Capelli, a. a. O. \S\ 23은 충분조건 $\bigl(a^{(k)}\mid \frac{\partial}{\partial a^{(i)}}\bigr)=0$, $(k>i,\ i=1,2,\ldots,\rho_1-1)$을 제시한다. 여기서 필요충분조건은 Wellstein, Von den Differentialgleichungen der projektiven Invarianten, Math. Ann. 67 (1909) \S\ 3에 따라 푸아송 괄호 과정을 적용하면 가장 간단히 얻어진다.}
\begin{equation}
 \left(a^{(i+1)}\middle|\frac{\partial}{\partial a^{(i)}}\right)=0;\ \ldots\ ;
 \left(k^{(j+1)}\middle|\frac{\partial}{\partial k^{(j)}}\right)=0.
 \quad (i=1,2,\ldots,\rho_1-1;\ j=1,2,\ldots,\rho_k-1)
\tag{19}
\end{equation}

% inherited: noether_paper04_ko_translation_001_20260804\targets\Noether_P04_Korean_T03_U14_UNCHECKED.tex
% sha256: F75C627C976334B74AC31BBC4C4E7DA5D7FFED9EE3F82F764177B3552DFCE508
이 합들이 모든 열을 명시적으로도 포함한다는 것을 보이자. $B^{\rho_2}\cdots K^{\rho_k}q_{n-\rho+1}=q_{n-\rho_1+1}$로 놓으면, (16.)과 (10.)에 따라 다음이 나온다.
\begin{align*}
\sum(-1)^{i+1}\pair{a^{(i)}}{x}(a^{(1)}\cdots a^{(i-1)}a^{(i+1)}\cdots a^{(\rho_1)}q_{n-\rho+1})
&=\pair{A^{\rho_1}}{xp_{\rho_1-1}}\\
&=(A_{n-\rho_1}xp_{\rho_1-1})=(-1)^{(\rho_1+1)(n+1)}\pair{q_{n-\rho_1+1}}{A_{n-\rho_1}x}\\
&=(-1)^{(\rho_1+1)(n+1)}\pair{B^{\rho_2}\cdots K^{\rho_k}q_{n-\rho+1}}{A_{n-\rho_1}x}.
\end{align*}
나머지 합들도 같은 방식으로 계산하고, 열들이 이미 무게지수로 서로 다르게 특징지어져 있으므로 나중에 $B^{\rho_2}=A^{\rho_2}$, $C^{\rho_3}=A^{\rho_3},\ldots$로 놓으면, 원하는 항등식을 얻는다.
\begin{equation}
\begin{aligned}
\pair{A^{\rho_1}A^{\rho_2}\cdots A^{\rho_k}}{xp_{\rho-1}}
&=\sum_{i=1}^{k}(-1)^{\delta_i}\pair{A^{\rho_1}\cdots A^{\rho_{i-1}}A^{\rho_{i+1}}\cdots A^{\rho_k}q_{n-\rho+1}}{A_{n-\rho_i}x},\\
\delta_i&=(\rho_1+\rho_2+\cdots+\rho_{i-1})\rho_i+(\rho_i+1)(n+1).
\end{aligned}
\tag{20}
\end{equation}
그리고 (17.)로부터 그 쌍대 항등식들을 얻는다.

% inherited: noether_paper04_ko_translation_001_20260804\targets\Noether_P04_Korean_T03_U15_UNCHECKED.tex
% sha256: 5039271D12EE05A91F907A02C0A43E0F1CAD07FCA54F8A868186422727DF5F94
\begin{equation}
\begin{aligned}
\pair{u q_{\sigma-1}}{A_{\sigma_1}A_{\sigma_2}\cdots A_{\sigma_k}}
&=\sum_{i=1}^{k}(-1)^{\eps_i}\pair{u A^{n-\sigma_i}}{p_{n-\sigma+1}A_{\sigma_1}\cdots A_{\sigma_{i-1}}A_{\sigma_{i+1}}\cdots A_{\sigma_k}},\\
\eps_i&=(\sigma_1+\sigma_2+\cdots+\sigma_{i-1})\sigma_i+(\sigma_i+1)(n+\sigma).
\end{aligned}
\tag{21}
\end{equation}
\emph{정리 II.} (20.)과 (21.)로 열 $A^{\rho_i},x,p$ 사이의, 또는 각각 $A_{\sigma_i},u,q$ 사이의 분해 불가능 항등식 전체가 소진된다. 동시에, 대입 (11.)을 하지 않고도 명시적 표현을 허용하는 분해 불가능 항등식 전체도 이들로 소진된다.

% inherited: noether_paper04_ko_translation_001_20260804\targets\Noether_P04_Korean_T03_U16_UNCHECKED.tex
% sha256: 9D63C88E326907CE38D0AE4A7C3970FF3DD57EE7D3565C1733E34BBCE5103993
정리의 첫째 부분은 (20.)과 (21.)을 유도할 때의 고찰에서 나오며, 둘째 부분은 다음과 같은 두 기호열 사이의 항등식이 존재할 수 없다는 데서 나온다.
\[
\pair{q_\rho A^\sigma}{B_{\rho+\sigma-\tau}p_\tau}
=\eps_1\pair{q_\rho B^{n-\rho-\sigma+\tau}}{A_{n-\sigma}p_\tau}
+\eps_2\pair{q_\rho q_{n-\tau}}{B_{\rho+\sigma-\tau}A_{n-\sigma}},
\]
여기서 $\rho\leqq\tau\leqq\sigma$이고 어느 열도 무게 $1\cdot(n-1)$을 갖지 않는다. ($\rho,\tau,\sigma$에 관한 다른 가정은 열의 표기를 서로 바꾸면 이 경우로 환원된다.) 예를 들어 이 불가능성은 (8.)을 전개하되 $q_\rho=lM^{\rho-1}$, $q_{n-\tau}=lN^{n-\tau-1}$로 특수화하여 얻는다. 이때 $l_1=1$, $M^{2,3,\ldots,\rho}=1$, $A^{\rho+1,\ldots,\rho+\sigma}=1$로 놓고 열 $l,M,A$의 나머지 모든 원소를 0으로 놓는 한편, $N$과 $B$는 임의로 둔다. 이제 임의의 열 사이의 항등식은 하나의 열 $p_\tau$를 $y^{(1)}y^{(2)}\cdots y^{(\tau)}$로 풀어 놓음으로써 (20.)으로부터 합성되는 항등식으로 바뀐다. 그러므로 (20.)을 기호열 두 개만을 가진 항등식들로부터 생성할 수 있음을 보이기만 하면 정리의 둘째 부분이 증명된다. 실제로
\begin{equation}
\pair{A^{\rho_1}A^{\rho_2}}{xp_{\rho-1}}
=\eps\pair{A^{\rho_1}q_{n-\rho+1}}{A_{n-\rho_2}x}
+\pair{A^{\rho_1}}{xp_{\rho_1-1}},
\quad \text{여기서 }p_{\rho_1-1}\sim A^{\rho_2}q_{n-\rho+1},
\tag{20a}
\end{equation}
에서 $A^{\rho_1}=B^{\sigma_1}B^{\sigma_2}$로 특수화한다. 즉,
\[
\pair{B^{\sigma_1}B^{\sigma_2}}{xp_{\rho_1-1}}
=\eps\pair{B^{\sigma_1}q_{n-\rho_1+1}}{B_{n-\sigma_2}x}
+\pair{B^{\sigma_1}}{xp_{\sigma_1-1}},
\quad \text{여기서 } p_{\sigma_1-1}\sim B^{\sigma_2}q_{n-\rho_1+1},
\]
로 특수화하면 세 기호열에 관한 하나의 항등식 (20.)을 얻는다. 그리고 $B^{\sigma_1}=C^{\tau_1}C^{\tau_2}$ 등으로 계속 특수화하면 일반 항등식 (20.)을 얻는다. 따라서 두 기호열과 두 임의의 변수열 사이의 항등식에는 명시적 표현이 존재하지 않는다는 사실로부터, 변수열들 가운데 열 $x$나 $u$가 하나도 없다면 기호열이 임의로 많이 있는 경우에도 명시적 표현이 존재하지 않는다는 결론이 따른다.

% inherited: noether_paper04_ko_translation_001_20260804\targets\Noether_P04_Korean_T04_U17_UNCHECKED.tex
% sha256: 48C5D717BF3DBDE03828A3CEF518B9984D1BC43BABA9BC6645940B09E1EDCB1D
\section*{\S\ 4. 분해 항등식.}
이제 임의의 기호열과 변수열 사이의 항등식 가운데 가장 일반적인 경우를 생각한다. 이 항등식의 종결식에는 대입 (11.)을 하지 않은 채로 하나의 열 $p_\tau$가 그 개별 열 $y^{(1)}\cdots y^{(\tau)}$로 분해되어 나타난다. 우리는 이러한 항등식을 ``분해 항등식''이라 부른다. 앞 절 끝의 주석에 이어서, 먼저 기호열이 두 개뿐인 경우에 이 항등식을 구한다. 즉, 식 $\pair{q_\rho A^\sigma}{B_{\rho+\sigma-\tau}p_\tau}$을 전개한다. 이때 다음과 같이 놓는다.
\begin{equation}
\begin{aligned}
q_{\rho-\alpha}^{(i_1\ldots i_\alpha)}&\sim p_{n-\rho}\,y^{(i_1)}\cdots y^{(i_\alpha)},\qquad
p_{\tau-\alpha}^{(i_1\ldots i_\alpha)}=y^{(i_{\alpha+1})}\cdots y^{(i_\tau)},\qquad
p_\tau=y^{(1)}\cdots y^{(\tau)}.
\end{aligned}
\tag{22}
\end{equation}
수 $\rho,\sigma,\tau$에 관해서는 네 경우를 구별해야 한다.
\[
\begin{array}{llll}
1)&\rho+\sigma\leqq n,& \tau\leqq\rho,& \tau\leqq\sigma;\\[2pt]
2)&\rho+\sigma\leqq n,& \tau\geqq\rho,& \tau\leqq\sigma;\\[2pt]
3)&\rho+\sigma\leqq n,& \tau\leqq\rho,& \tau>\sigma;\\[2pt]
4)&\rho+\sigma\leqq n,& \tau\geqq\rho,& \tau\geqq\sigma,
\end{array}
\]
이 가운데 첫째 경우를 택하여 다룬 뒤 나머지 경우를 짧게 특징짓겠다. 열 $y^{(1)}y^{(2)}\cdots y^{(\tau)}B_{\rho+\sigma-\tau}$를 $y^{(1)}$과 $y^{(2)}\cdots y^{(\tau)}B_{\rho+\sigma-\tau}$로 분해하면, (20.), (10.), (9.)에 따라 다음이 나온다.

% inherited: noether_paper04_ko_translation_001_20260804\targets\Noether_P04_Korean_T04_U18_UNCHECKED.tex
% sha256: 7A5EECB9CE6AF71B372A293ED23113913CFCB8CB99454B16EFD1B181E4B1B425
\begin{equation}
\begin{aligned}
\pair{q_\rho A^\sigma}{y^{(1)}y^{(2)}\cdots y^{(\tau)}B_{\rho+\sigma-\tau}}
&=\pair{q_{\rho-1}^{(1)}A^\sigma}{y^{(2)}\cdots y^{(\tau)}B_{\rho+\sigma-\tau}}\\
&\quad+(-1)^\rho\pair{q_\rho[A_{n-\sigma}y^{(1)}]^{\sigma-1}}{y^{(2)}\cdots y^{(\tau)}B_{\rho+\sigma-\tau}}.
\end{aligned}
\tag{23}
\end{equation}
(23.)의 마지막 항은 원래 식과 정확히 같은 형태를 가진다. 다만 $\sigma$가 $\sigma-1$로, $\tau$가 $\tau-1$로 바뀌었을 뿐이다. 따라서 이 항에도 다시 등식 (23.)을 적용할 수 있다. 그러므로 $\tau\leqq\sigma$이므로 연산 (23.)을 $\tau$번 적용하고 (10.)을 고려하면 다음 관계에 이른다.
\begin{equation}
\begin{aligned}
\pair{q_\rho A^\sigma}{p_\tau B_{\rho+\sigma-\tau}}
&=(-1)^{\rho\sigma+(\sigma+\tau)(n+1)}
  \pair{q_\rho B^{n-\rho-\sigma+\tau}}{A_{n-\sigma}p_\tau}\\
&\quad+\sum_{i_1=1}^{\tau}(-1)^{\rho(i_1-1)}
 \pair{q_{\rho-1}^{(i_1)}[A_{n-\sigma}y^{(1)}\cdots y^{(i_1-1)}]^{\sigma-i_1+1}}
 {y^{(i_1+1)}\cdots y^{(\tau)}B_{\rho+\sigma-\tau}}.
\end{aligned}
\tag{24}
\end{equation}

% inherited: noether_paper04_ko_translation_001_20260804\targets\Noether_P04_Korean_T04_U19_UNCHECKED.tex
% sha256: 69DBA358AC9882ED3CD7935A793986281858E6F8E614733E70B8D2D689FB7E1B
합 기호 아래의 각 항은 다시 원래 식의 유형을 가진다. $\rho$는 $\rho-1$로, $\sigma$는 $\sigma-i_1+1$로, $\tau$는 $\tau-i_1$로 바뀌었으므로 조건 1)은 오히려 더 강하게 충족된다. 따라서 $\tau\leqq\rho$이므로 연산 (24.)을 $\tau$번 적용할 수 있고, 원하는 분해 항등식을 얻는다.
\begin{equation}
\begin{aligned}
\pair{q_\rho A^\sigma}{p_\tau B_{\rho+\sigma-\tau}}
&=\eps_{i_0}\pair{q_\rho B^{n-\rho-\sigma+\tau}}{A_{n-\sigma}p_\tau}\\
&\quad+\sum_{i_1=1}^{\tau}\eps_{i_1}
 \pair{q_{\rho-1}^{(i_1)}B^{n-\rho-\sigma+\tau}}{A_{n-\sigma}p_{\tau-1}^{(i_1)}}\\
&\quad+\sum_{i_2=i_1+1}^{\tau}\eps_{i_2}
 \pair{q_{\rho-2}^{(i_1i_2)}B^{n-\rho-\sigma+\tau}}{A_{n-\sigma}p_{\tau-2}^{(i_1i_2)}}+\cdots\\
&\quad+\eps_{i_\tau}\pair{q_{\rho-\tau}^{(i_1\ldots i_\tau)}B^{n-\rho-\sigma+\tau}}{A_{n-\sigma}}\\
&=\sum_{\alpha=0}^{\tau}\sum_{i_\alpha=i_{\alpha-1}+1}^{\tau}\eps_{i_\alpha}
 \pair{q_{\rho-\alpha}^{(i_1\ldots i_\alpha)}B^{n-\rho-\sigma+\tau}}{A_{n-\sigma}p_{\tau-\alpha}^{(i_1\ldots i_\alpha)}}.
\end{aligned}
\tag{25}
\end{equation}
부호 $\eps$에 대해서는 (24.)로부터 다음을 얻는다.
\begin{equation}
\begin{aligned}
\eps_{i_\alpha}&=(-1)^{\delta_{i_\alpha}},\qquad
\delta_{i_\alpha}=\sum_{\beta=1}^{\alpha}(\rho-\beta+1)(i_{\beta-1}+i_\beta+1)\\
&\quad +(\rho-\alpha)(\sigma+i_\alpha+\alpha)+(\sigma+\tau+\alpha)(n+1),
\end{aligned}
\tag{26}
\end{equation}
여기서 $i_0=0$으로 놓아야 한다. (25.)의 합 기호는 다음과 같이 이해한다. $i_1<i_2<\cdots<i_{\alpha-1}$인 각 조합 $i_1i_2\ldots i_{\alpha-1}$에 대하여 $i_{\alpha-1}<i_\alpha\leqq\tau$를 만족하는 $i_\alpha$의 모든 값을 덧붙인다. (25.)의 개별 합들은 (26.)을 고려하면 쉽게 확인할 수 있듯이 (19.)와 유사한 다음 미분방정식을 만족한다.
\begin{equation}
 \left(\frac{\partial}{\partial y^{(i)}}\middle|y^{(i+1)}\right)=0,
 \qquad (i=1,2,\ldots,\tau-1),
\tag{27}
\end{equation}
따라서 이 합들은 열 $p_\tau$에만 의존한다. 그러므로 이 분해 항등식은 열 $A,B,q,p$ 사이의 분해 불가능 항등식이다. 명시적 표현은 다음 절에서 다시 다룬다.

% inherited: noether_paper04_ko_translation_001_20260804\targets\Noether_P04_Korean_T04_U20_UNCHECKED.tex
% sha256: 12292B0A0178D9DBA60498C8A84143624B6E8716F18FD79E5B9557D3B78B316B
경우 2)에서도 연산 (23.)은 $\tau$번 적용할 수 있지만, 연산 (24.)는 $\rho$번째 뒤에 중단된다. 따라서 (25.)의 종결식에서는 첫째 합 기호를 $\sum_{\alpha=0}^{\rho}$로 바꾸어야 한다. 경우 3)과 4)에서는 연산 (23.)을 $\sigma$번만 실행할 수 있다. 이때 (24.) 대신 다음 관계가 들어간다.
\begin{equation}
\begin{aligned}
\pair{q_\rho A^\sigma}{p_\tau B_{\rho+\sigma-\tau}}
&=(-1)^{\rho\sigma}
 \pair{q_{\rho-(\tau-\sigma)}^{(\sigma+1\ldots\tau)}B^{n-\rho+\tau-\sigma}}{A_{n-\sigma}p_{\tau-(\tau-\sigma)}^{(\sigma+1\ldots\tau)}}\\
&\quad+\sum_{i_1=1}^{\sigma}(-1)^{\rho(i_1-1)}
 \pair{q_{\rho-1}^{(i_1)}[A_{n-\sigma}y^{(1)}\cdots y^{(i_1-1)}]^{\sigma-i_1+1}}
 {y^{(i_1+1)}\cdots y^{(\tau)}B_{\rho+\sigma-\tau}};
\end{aligned}
\tag{28}
\end{equation}
(28.)을 $(\tau-\sigma)$번 반복하면, 그동안 합 기호는 $\sum_{i_\beta=i_{\beta-1}+1}^{\sigma+\beta-1}$로 바뀌며, (25.)에서 지수 $\alpha=\tau-\sigma$에 해당하는 개별 합의 모든 항을 각각 알맞은 부호와 함께 얻는다. 이는 (28.)의 합 기호 아래 항들에 (23.)을 $(\sigma-i_1+1)$번 적용하면 알 수 있다. 그 뒤에는 경우 3)과 4)가 경우 1)과 2)로 넘어간다. 실제로 $\sigma,\tau$에 대응하는 수들은 $\sigma'=\sigma-i_{\tau-\sigma}+\tau-\sigma$, $\tau'=\tau-i_{\tau-\sigma}=\sigma'$가 되며, 절차를 계속하면 $\tau'<\sigma'$가 된다. 그러므로 (25.)의 종결식에서는 첫째 합 기호를 $\sum_{\alpha=\tau-\sigma}^{\tau,\rho}$로 바꾸어야 한다.

% inherited: noether_paper04_ko_translation_001_20260804\targets\Noether_P04_Korean_T04_U21_UNCHECKED.tex
% sha256: FE69137C3446CCD17F66D9FC33ECA71E3809BEDD9BFE6359407682ABD8B348C9
(21.)로부터 이와 유사하게 쌍대적인 분해 항등식을 얻는다. 여기서는 하나의 열 $q_\rho$가 개별 열 $v^{(1)},\ldots,v^{(\rho)}$로 분해된다. 다음과 같이 놓자.
\begin{equation}
 \dot p_{\tau-\beta}^{(i_1\ldots i_\beta)}\sim v^{(i_1)}\cdots v^{(i_\beta)}q_{n-\tau};\qquad
 \dot q_{\rho-\beta}^{(i_1\ldots i_\beta)}=v^{(i_{\beta+1})}\cdots v^{(i_\rho)};\qquad
 q_\rho=v^{(1)}\cdots v^{(\rho)},
\tag{29}
\end{equation}
그러면 다음이 나온다.
\begin{equation}
 \pair{q_\rho A^\sigma}{p_\tau B_{\rho+\sigma-\tau}}
 =\sum_{\substack{0,\tau-\sigma}}^{\substack{\tau,\rho}}
   \sum_{i_\beta=i_{\beta-1}+1}^{\rho}\eps
\pair{\dot q_{\rho-\beta}^{(i_1\ldots i_\beta)}B^{n-\rho-\sigma+\tau}}{A_{n-\sigma}\dot p_{\tau-\beta}^{(i_1\ldots i_\beta)}} ,
\tag{30}
\end{equation}
여기서 합 지수 $\beta$는 아래의 두 값 가운데 큰 값부터 위의 두 값 가운데 작은 값까지 변한다. (25.)와 (30.)의 한 개별 합에 속하는 모든 항은 하나의 형식을 다음 편극 과정들로부터 얻은 편극들이다.
\[
\left(\frac{\partial}{\partial p_{\tau-\alpha}}\middle|p_{\tau-\alpha}^{(i_1\ldots i_\alpha)}\right),\qquad
\left(q_{\rho-\alpha}^{(i_1\ldots i_\alpha)}\middle|\frac{\partial}{\partial q_{\rho-\alpha}}\right).
\]
이를 표현하기 위하여 상수와 곱한 이러한 편극 연산들의 결합식을 $\Pi$로 나타내기로 한다. 그러면 (25.)와 (30.) 대신 다음 관계를 얻는다.
\begin{equation}
 \pair{q_\rho A^\sigma}{p_\tau B_{\rho+\sigma-\tau}}
 =\sum_{\substack{0,\tau-\sigma}}^{\substack{\tau,\rho}}
 \Pi\pair{q_{\rho-\alpha}B^{n-\rho-\sigma+\tau}}{A_{n-\sigma}p_{\tau-\alpha}}.
\tag{31}
\end{equation}

% inherited: noether_paper04_ko_translation_001_20260804\targets\Noether_P04_Korean_T05_U22_UNCHECKED.tex
% sha256: 372CAE0ACCFF3DC0AA4EC8313A6F09D67C8911F634D7A27287BCC1690F5F2146
\section*{\S\ 5. 명시적 형태의 분해 항등식.}
경우 4)의 (25.)에서 $\tau$를 $\sigma+\rho$로 특수화하면 다음이 나온다.
\begin{equation}
 \pair{q_\rho A^\sigma}{y^{(1)}\cdots y^{(\rho+\sigma)}}
 =\sum_{i_\rho=i_{\rho-1}+1}^{\tau}
 \eps_{i_\rho}\pair{q_\rho}{y^{(i_1)}\cdots y^{(i_\rho)}}
 \pair{A^\sigma}{y^{(i_{\rho+1})}\cdots y^{(i_{\rho+\sigma})}},
\tag{32}
\end{equation}
여기서 $\eps_{i_\rho}=(-1)^{\delta_{i_\rho}}$, $\delta_{i_\rho}=\rho(\rho+1)/2+i_1+\cdots+i_\rho$가 된다. 이것은 (16.)과 (17.)보다 더 일반적인 형태의 행렬 곱정리이며, 곱행렬의 행렬식
\[
 \sum\pm(v^{(1)}y^{(1)})\cdots(v^{(\rho)}y^{(\rho)})(a^{(1)}y^{(\rho+1)})\cdots(a^{(\sigma)}y^{(\rho+\sigma)})
\]
을 라플라스 전개하여 얻는다. 따라서 더 일반적인 곱정리는 항등식 (20.), 또는 각각 (21.)로부터 합성되며, 이로써 특수한 형태 (16.), (17.)을 따로 골라낸 이유가 설명된다.

% inherited: noether_paper04_ko_translation_001_20260804\targets\Noether_P04_Korean_T05_U23_UNCHECKED.tex
% sha256: 7F4659F3777765F55BBA59AC1C4F04B4EA229E0A9E394512510E227BB4296C3B
이제 관계 (32.)는 분해 항등식을 명시적으로 표현하는 수단을 준다. 이를 위하여 (31.)에 나타나는 형식들을 기본형으로 본다. 즉, 대입 (11.)을 실행한다. (8.)에 따르면 새로 도입한 열 $R$은 열 $A,B$ 자체로 표현할 수 있고 그 부분열 $a^{(1)}a^{(2)}\ldots$, $b^{(1)}b^{(2)}\ldots$로 표현할 필요가 없으므로, 이 대입은 열 $A,B$에 관한 명시적 표현을 바꾸지 않는다. 따라서
\begin{equation}
 \pair{q_{\rho-\alpha}B^{n-\rho-\sigma+\tau}}{A_{n-\sigma}p_{\tau-\alpha}}
 =\bigl(R_{\rho-\alpha}p_{n-\rho+\alpha}\bigr)
  \pair{R^{\tau-\alpha}}{p_{\tau-\alpha}},
\tag{33}
\end{equation}
라 하면, (25.)에서 지수 $\alpha$에 해당하는 개별 합에 대하여 다음이 나온다.
\begin{equation}
\begin{aligned}
&\sum \eps_{i_\alpha}\bigl(R_{\rho-\alpha}p_{n-\rho}y^{(i_1)}\cdots y^{(i_\alpha)}\bigr)
 \pair{R^{\tau-\alpha}}{y^{(i_{\alpha+1})}\cdots y^{(i_\tau)}}\\
&\quad=\sum \eps_{i_\alpha}\bigl([R_{\rho-\alpha}p_{n-\rho}]^\alpha y^{(i_1)}\cdots y^{(i_\alpha)}\bigr)
 \pair{R^{\tau-\alpha}}{y^{(i_{\alpha+1})}\cdots y^{(i_\tau)}}\\
&\quad=\eps\pair{[R_{\rho-\alpha}p_{n-\rho}]^\alpha R^{\tau-\alpha}}{p_\tau}
 =\eps\pair{R^{\tau-\alpha}q_{n-\tau}}{R_{\rho-\alpha}p_{n-\rho}}.
\end{aligned}
\tag{34}
\end{equation}
이로써 실제로 명시적인 종결식이 주어진다. $q_\rho$와 $p_\tau$가 들어가는 대칭적인 형태는 (30.)도 같은 종결식에 이른다는 것을 보여 준다. 따라서 이 종결식은 처음의 분해 방식에 의존하지 않는다. (25.)와 (30.) 사이의 차이는 열 $y$와 $v$가 독립적인 의미를 가질 때에만 존재한다. 이때는 우리의 정의에 따른 분해 항등식이 하나의 분해 가능한 항등식으로 바뀐다. 더 많은 기호열을 가진 항등식으로 넘어가려면 \S\ 3 끝에서처럼 한 열 $q_\rho$를 $q_{\rho_1}C^{\sigma_1}$로, 또는 각각 $p_\tau$를 $p_{\tau_1}C_{\tau-\tau_1}$로 풀어 놓기만 하면 된다. 그러면 생기는 식
\begin{equation}
 \pair{q_{\rho_1}C^{\sigma_1}}{[R^{\tau-\alpha}q_{n-\tau}]_\lambda R_{\rho_1+\sigma_1-\alpha-\lambda}},
 \qquad
 \pair{[R_{\rho-\alpha}p_{n-\rho}]^\lambda R^{\tau-\alpha}}{p_{\tau_1}C_{\tau-\tau_1}},
\tag{35}
\end{equation}
은 두 기호열의 경우에 나타나는 유형과 정확히 같다. 따라서 (33.)에 해당하는 대입을 하면 다시 명시적 표현을 허용하며, 이 과정을 반복하면 임의로 많은 열에 대해서도 명시적 표현을 얻는다. 또한 (25.)와 각각 (30.)의 첫째 합과 마지막 합은 (33.)을 적용하지 않고 공식 (32.)만으로 명시적 표현을 주거나, 아예 모든 열을 명시적으로 포함하는 단 하나의 항으로 줄어든다. (32.)에서 $q_\rho$를 $q_{\rho_1}C^{\sigma_1}$ 등으로 풀어 얻는 관계는 가장 일반적인 형태의 곱정리로 볼 수 있다.

% inherited: noether_paper04_ko_translation_001_20260804\targets\Noether_P04_Korean_T05_U24_UNCHECKED.tex
% sha256: E1D6918969032A8D6C586409B474AE99F6B76D89E9E7C576B106CA1FA3AD1AF5
요약하면 우리는 다음을 보일 것이다.

\emph{정리 III:} 항등식
\begin{equation}
 \pair{q_\rho A^\sigma}{p_\tau B_{\rho+\sigma-\tau}}
 =\sum_{\substack{0,\tau-\sigma}}^{\substack{\tau,\rho}}
 \eps\pair{R^{\tau-\alpha}q_{n-\tau}}{R_{\rho-\alpha}p_{n-\rho}},
\tag{36}
\end{equation}
에서 열 $R$은 (33.)으로 정해진다. 이 항등식들과, 여기서 $q_\rho,p_\tau$를 $q_{\rho_1}C^{\sigma_1}$ 등으로 풀어 얻는 항등식들로 분해 불가능 항등식 전체가 소진된다.

실제로 이렇게 생기는 항등식들은 그 종류 가운데 가장 일반적이다. (20.), (21.)에서는 아직 존재하던 무게와 개수에 관한 제한을 이제는 개별 열들이 더 이상 받지 않기 때문이다. 또한 이 열들 사이에는 다른 항등식도 존재하지 않는다. 한 열 $p_\tau$를 $y^{(1)}\cdots y^{(\tau)}$로 풀면 가장 일반적인 항등식은 (20.)으로부터, 따라서 \S\ 3 끝에 따르면 (20a.)로부터 합성된다. 그리고 (16.) 뒤에서 논한 합성은 정확히 \S\ 4에서 실행한 합성이다. 임의로 많은 열로 넘어가는 것은, (20a.)의 논의가 보여 주듯이, $q_\rho$를 $q_{\rho_1}C^{\sigma_1}$ 등으로 풀어 얻은 식들에 언제나 같은 연산을 적용함으로써 이루어진다. 이로부터 (20a.) 대신 (20.)에서 직접 넘어가더라도 새로운 것은 생기지 않는다는 결론도 따른다. 이는 한 번은 (25.)에서 열 $q_\rho$를 특수화하고, 다른 한 번은 \S\ 4의 방법으로 (20.)에서 넘어가서 계산하면 쉽게 확인할 수 있다. 두 경우 모두 같은 합들이 생기며, 일반 곱정리(위 참조)를 적용하면 (36.)에서 얻는 항등식으로 바뀐다. 항등식 (20.), (21.)은 특수한 경우 $\tau=1$, $\rho=1$에 해당한다. 이때는 각각 두 개의 개별 합만 나타나므로, 위 주석에 맞게 대입 (33.) 없이 명시적 표현을 얻으며, 이는 앞서 얻은 결과와 일치한다.

% inherited: noether_paper04_ko_translation_001_20260804\targets\Noether_P04_Korean_T05_U25_UNCHECKED.tex
% sha256: 432EB596F5FA17A35B64916810C3A26A55499EF0DBD3EB84A6559FC0778A8A88
끝으로 분해 항등식의 두 특수한 경우를 짧게 언급하자. (31.)에서 $\tau=\sigma$, $\rho=n-\sigma$로 놓고, 열 $p_\tau$를 $p'_\sigma\sim q'_{n-\sigma}$로 나타내면 다음이 나온다.
\begin{equation}
\begin{aligned}
\pair{A^\sigma}{p_\sigma}\pair{B^\sigma}{p'_\sigma}
-\pair{A^\sigma}{p'_\sigma}\pair{B^\sigma}{p_\sigma}
=\sum_{\alpha=1}^{\sigma,n-\sigma}
\Pi\pair{q_{n-\sigma-\alpha}B^\sigma}{A_{n-\sigma}p_{\sigma-\alpha}}.
\end{aligned}
\tag{37}
\end{equation}
이것은 서로 동변인 두 열 $p_\sigma,p'_\sigma$를 가진 한 형식을 기본 공변식들로 전개하는 Clebsch 공식이다.\footnote{Clebsch, a. a. O., S. 25--32. Clebsch는 오른쪽의 개별 항들이 열 $A,B$에만 의존한다는 것을 보였다. 그의 식 $\Phi_h$에 (32.)를 적용하면 명시적 표현을 얻을 수 있을 것이다.} 형식 $\pair{A^\sigma}{p_\sigma}^m\pair{B^\sigma}{p'_\sigma}^n$에 속하는 기본 공변식들은 다음과 같다.
\begin{equation}
 \left\{\sum_{\alpha=1}^{\sigma,n-\sigma}
 \pair{q_{n-\sigma-\alpha}B^\sigma}{A_{n-\sigma}p_{\sigma-\alpha}}
 \right\}^{\rho}
 \pair{A^\sigma}{p_\sigma}^{m-\rho}\pair{B^\sigma}{p'_\sigma}^{n-\rho},
 \qquad (\rho=0,1,\ldots,m,n).
\tag{38}
\end{equation}
이를 간단히 말하면, 기본형으로부터 ``결손 $\rho$의 동변 수축''(\S\ 2 참조)을 하여 얻는다.

% inherited: noether_paper04_ko_translation_001_20260804\targets\Noether_P04_Korean_T05_U26_UNCHECKED.tex
% sha256: 3EF6BAEF876C909D76FF83A1324BCE3640F2DBA410A42DA4908FB61CEE8B1E6F
둘째로 $\tau=\sigma-\lambda$, $\rho=n-\sigma-\lambda$, $B^\sigma=A^\sigma$로 놓으면 다음이 나온다.
\begin{equation}
\begin{aligned}
\pair{q_{n-\sigma-\lambda}A^\sigma}{A_{n-\sigma}p_{\sigma-\lambda}}
&=(-1)^\lambda\pair{q_{n-\sigma-\lambda}A^\sigma}{A_{n-\sigma}p_{\sigma-\lambda}}\\
&\quad+(-1)^{(n-\sigma)(\sigma-\lambda)}
 \sum_{\alpha=1}^{\sigma-\lambda,n-\sigma-\lambda}
 \Pi\pair{q_{n-\sigma-\lambda-\alpha}A^\sigma}{A_{n-\sigma}p_{\sigma-\lambda}}.
\end{aligned}
\tag{39}
\end{equation}
따라서 \S\ 2에서 언급했듯이, 기호열을 특징짓는 조건 (12.) 가운데 결손 $\lambda$가 홀수인 조건은 더 높은 결손의 조건들로부터 따른다. 선형형식 $\pair{A^\sigma}{p_\sigma}=\pair{B^\sigma}{p_\sigma}$에 대해서는 (39.)로부터, 계수에 관해 이차인 그 기약 불변 구성식들이 짝수 결손의 구성식들로 환원된다는 사실을 얻는다. 이는 이항 및 삼항 영역에서 하나의 선형형식이 불변 구성식을 갖지 않는다는 알려진 사실과 일치한다.

일반 항등식들은 처음에 개별 열들이 조건 (12.)를 만족한다는 가정 아래 유도되었다는 점을 덧붙인다. 그러나 이 개별 열들은 항등식에 선형으로만 들어가므로, 이 항등식들은 조건 (12.)를 만족하지 않는 더 일반적인 열들에 대해서도 성립한다.

% inherited: noether_paper04_ko_translation_001_20260804\targets\Noether_P04_Korean_T06_U27_UNCHECKED.tex
% sha256: 5ACE78C48EEC3E2D77141EA8177F2737D9D1ED3392E465DDB080EDC0CF654760
\section*{\S\ 6. 불변 구성식의 명시적 표현.\\수축 과정과 미분 과정.}
앞 절들의 결과를 이용하면, \S\ 2에서 진술한 기호적 표현 가능성에 관한 정리(첫째 기본정리)에 명시적 표현 가능성을 덧붙여 완결할 수 있다. 즉, 다음을 보인다.

\emph{정리 IV:} ``임의의 기본형들의 모든 불변 구성식은 행렬곱으로 명시적으로 표현할 수 있다. 다시 말해, 종결식에는 변수열 $p_\rho$ 이외에 원래 형식들의 열 $S^\sigma,\ldots,T^\tau$나, 이들로부터 대입 (9.) 또는 (11.)로 유도한 열 $P,Q,R$만이 들어간다.''

증명을 위하여 한 불변 구성식이 변수열뿐 아니라 열 $S^\sigma,\ldots,T^\tau$에도 의존한다고 하자. 이 열들은 행렬식에 의한 표현이 가능하도록 개별 열들로 충분히 세분되어 있다고 하자. 즉, $S^\sigma=S^{\sigma_1}\cdots S^{\sigma_i},\ldots,T^\tau=T^{\tau_1}\cdots T^{\tau_k}$이다. 열들은 그 자체로는 임의이지만 확정해 둔 순서 $S^\sigma,\ldots,T^\tau$로 배열되어 있다고 하자. 이제 첫째 전체 열 $S^\sigma=S^{\sigma_1}\cdots S^{\sigma_i}$에, 그 부분열 몇 개에만이 아니라 전체에 의존하는 행렬식 결합식 하나를 택한다. 그런데 이 의존성은 \S\ 3--\S\ 5에 따르면 일반 항등식들의 귀결이다. 이제 한 변수열 $p_\rho$를 열 $y,v$로 풀거나, 뒤에 오는 열 $T$ 가운데 하나를 개별 열 $t$, 또는 각각 $\tau$로 풀면, 가장 일반적인 항등식들은 (20.), (21.)로부터 합성된다. 그런데 이 항등식들은 언제나 열 $S^\sigma=S^{\sigma_1}\cdots S^{\sigma_i}$를 명시적으로 포함하는 식, 즉 (20.), (21.)의 왼쪽으로 이끈다. 실제로 (19.)를 보면 오른쪽의 완전한 합 전체만이 --- (20a.)의 한 항에 해당하는 그 일부가 아니라 --- $S^\sigma$에만 의존한다는 성질을 가짐을 알 수 있다. 절차를 계속하는 동안 한 번 명시적으로 들어간 모든 열은 계속 명시적으로 남는다. 항등식을 적용하면서 요구될 수 있는 것은 여러 열을 하나로 묶는 것, 즉 대입 (9.)뿐이다. 마지막으로 단 하나의 열 $T^\tau$만 명시적 형태로 바꾸면 되는 단계에서는 두 경우가 가능하다. 자유로운 변수열, 곧 대입 (9.)로 다른 열과 결합되지 않은 변수열이 아직 남아 있을 수 있다. 이를 부분열 $y$ 또는 $v$로 분해하면 절차를 끝낼 수 있고, (31.)에서처럼 모든 열을 명시적으로 포함하는 하나 이상의 형식의 편극들이 생긴다. 자유로운 변수열이 더 이상 없다면, 개별 열 $t$ 또는 $\tau$를 포함하면서 전체 열 $T^\tau$에 의존하는 식들의 총체에서, 그 생성 방식에 따라 분해 항등식의 항들을 알아볼 수 있다. 그런데 이 항들은 \S\ 5에 따라 대입 (11.)을 적용하면 언제나 모든 열에 관한 명시적 표현을 허용한다. 이로써 위 정리는 일반적으로 증명되었다. 기호열이 두 개만 나타날 때에는 대입 (11.)이 전혀 일어나지 않는다는 점도 덧붙인다. 실제로 $\sigma,n-\sigma,\tau,n-\tau<n$이므로, 변수열이 전혀 들어가지 않을 수 있는 것은 두 기호열이 하나의 행렬식에 결합되어 이미 명시적으로 나타나는 경우뿐이다. 반대로 변수열이 들어가면 (34.)와 (33.)에 따라 대입 (11.)은 불필요해진다.

% inherited: noether_paper04_ko_translation_001_20260804\targets\Noether_P04_Korean_T06_U28_UNCHECKED.tex
% sha256: 5AF078024CBEC085E7C7AF45B655741157C8814EFC660B1594D69844E37FF79C
방금 말한 사실에 따르면 모든 불변 구성식을 생성하기 위해서는 기호열들에 변수열을 덧붙여 가능한 모든 행렬곱으로 결합하는 것으로 충분하다. 가장 일반적인 행렬곱은 다음 형태이다.
\begin{equation}
 \pair{P^{\mu_1}\cdots P^{\mu_i}}{Q_{\nu_1}\cdots Q_{\nu_k}},\qquad \sum\mu=\sum\nu,
\tag{40}
\end{equation}
여기서 $P$와 $Q$는 원래 형식들의 열일 수도 있고, 원래 열들에 대입 (9.) 또는 (11.)을 차례로 적용하여 생긴 열일 수도 있으며, 마지막으로 변수열일 수도 있다. 그런데 식 (40.)은 변수열을 이용하면서 매번 두 기호열을 하나의 행렬곱으로 차례로 결합하여 생성할 수 있다. 실제로
\[
 \pair{q_{n-\mu-\lambda}P^\mu}{Q_{\nu_1}p_{n-\nu_1-\lambda}}
 =\bigl(R_\lambda Q_{\nu_1}p_{n-\nu_1-\lambda}\bigr)
\]
에 열 $Q_{\nu_2}$를 결합하면 식
\[
 \pair{[R_\lambda Q_{\nu_1}]^{n-\lambda-\nu_1}}{Q_{\nu_2}p_{n-\nu_1-\nu_2-\lambda}}
 =\pair{q_{n-\mu-\lambda}P^\mu}{Q_{\nu_1}Q_{\nu_2}p_{n-\nu_1-\nu_2-\lambda}}
\]
을 얻는다. 따라서 절차를 계속하면 식 (40.)을 얻는다. 이때 $P$와 $Q$가 원래 형식들의 열이 아니라면, (9.)와 (11.)을 방금 보인 사실과 함께 적용하여 이 열들도 매번 두 기호열을 결합하는 절차를 차례로 거쳐 생겼음을 알 수 있다. 따라서 다음이 성립한다.

\emph{정리 V:} 모든 불변 구성식을 생성하는 과정인 수축 과정은, 변수열을 이용하면서 매번 두 기호열을 하나의 행렬곱으로 차례로 결합하는 데 있다.

% inherited: noether_paper04_ko_translation_001_20260804\targets\Noether_P04_Korean_T06_U29_UNCHECKED.tex
% sha256: 51D5EF319190C3D155F1049CDBE7958142436585EB7ACA0C0EC89CB5F978763A
$\sigma\geqq\tau$인 두 기호열 $S^\sigma,T^\tau$로부터 다음 행렬곱들을 만들 수 있다.
\begin{align}
&\pair{q_{n-\sigma-\lambda}S^\sigma}{T_{n-\tau}p_{\tau-\lambda}}, &&(\lambda=1,2,\ldots,\tau,n-\sigma),\tag{41}\\
&\pair{q_{n-\tau-\mu}T^\tau}{S_{n-\sigma}p_{\sigma-\mu}}, &&(\mu=\sigma-\tau+\lambda),\tag{42}\\
&\pair{q_{n-\tau-\nu}T^\tau}{S_{n-\sigma}p_{\sigma-\nu}}, &&(\nu=1,2,\ldots,\sigma-\tau).
\tag{43}
\end{align}
그러나 (31.)에서 (42.)와 (43.)에 대하여 다음 값들을 얻는다.
\begin{align}
&\pair{q_{n-\sigma-\lambda}S^\sigma}{T_{n-\tau}p_{\tau-\lambda}}
 +\sum_{\alpha}\Pi\pair{q_{n-\sigma-\lambda-\alpha}S^\sigma}{T_{n-\tau}p_{\tau-\lambda-\alpha}},\tag{42a}\\
&\Pi(q_{n-\sigma}S^\sigma)(T_{n-\tau}p_\tau)
 +\sum_{\beta}\Pi\pair{q_{n-\sigma-\beta}S^\sigma}{T_{n-\tau}p_{\tau-\beta}}.
\tag{43a}
\end{align}
따라서 형식 (42.)는 (41.) 및 (41.)의 편극들의 선형식으로 나타낼 수 있고, 형식 (43.)는 기본형 및 형식 (41.)의 편극들로 나타낼 수 있다. 이와 마찬가지로 (31.)에서 형식 (41.)이 (42.)와 (42.)의 편극들에 선형으로 의존함을 얻는다. 그러므로 Clebsch의 정의에 따르면\footnote{Clebsch, a. a. O. \S\ 7.} 형식 (42.)는 형식 (41.)와 동치인 반면, (43.)는 기본형으로 되돌아간다. 따라서 생성 과정은 (41.) 또는 (42.)로 동등하게 주어진다. 우리는 수 $\lambda$에 관해 더 단순한 과정 (41.)을 수축 과정으로 정의하겠다. 수축의 결손인 수 $\lambda$(\S\ 2 참조)는 행렬식의 곱에 도달하기까지 부족한 열의 개수를 나타낸다. 정확히는 주어진 수축에서 최대 개수의 열을 포함하는 그 행렬곱에서 부족한 열의 수이다. $\sigma\geqq\tau$일 때 $\lambda$는 두 값 $\tau,n-\sigma$ 가운데 작은 값까지만 변하므로 다음이 성립한다.
\begin{equation}
 \lambda\leqq\frac n2,\ n\text{ 짝수};\qquad
 \lambda\leqq\frac{n-1}{2},\ n\text{ 홀수}.
\tag{44}
\end{equation}

% inherited: noether_paper04_ko_translation_001_20260804\targets\Noether_P04_Korean_T06_U30_UNCHECKED.tex
% sha256: 5052FF31E086BA06D5ECFE4D039A171E3F55E2055607462356B3F471181DE513
추가 수축으로 열 $p,q$가 기호열로 바뀐 경우에도 (42.)와 (43.)을 (41.)로 환원하는 것은 여전히 성립한다. 실제로 (36.)에 따라
\[
 \pair{A^{n-\tau-\varkappa}T^\tau}{S_{n-\sigma}B_{\sigma-\varkappa}}
 =\sum_{\alpha=0,\varkappa-\sigma+\tau}^{\tau,n-\sigma}
 \eps\pair{R^{\tau-\alpha}B^{n-\sigma+\varkappa}}{A_{\tau+\varkappa}R_{n-\sigma-\alpha}},
\]
이고, 여기서 열 $R$은 (33.)에 따라 $S$와 $T$에서 생긴다. 따라서 기호열만을 포함하는 이 형식들은, $\sigma-\tau\gtreqless 2\varkappa$인지에 따라 각각 $\pair{q_\tau A^{n-\tau-\varkappa}}{B_{\sigma-\varkappa}p_{n-\sigma}}$ 또는 $\pair{q_{\tau-\alpha}B^{n-\sigma+\varkappa}}{A_{\tau+\varkappa}p_{n-\sigma-\alpha}}$를 기본형 및 형식 (41.)과 수축하여 얻는다.

% inherited: noether_paper04_ko_translation_001_20260804\targets\Noether_P04_Korean_T06_U31_UNCHECKED.tex
% sha256: 74141212EE4CFAE1B645DB35E79A730DAE7CF2BAEF0782C46604F2D863C5A185
기호적 수축 과정에서 불변 미분 과정을 얻는 방법은 잘 알려져 있다. 기호열 $S^\sigma,T_{n-\tau}$를 미분기호의 열 $\frac{\partial}{\partial p_\sigma}$, $\frac{\partial}{\partial q_{n-\tau}}$로 바꾸기만 하면 된다. 이때 각 곱 $\frac{\partial}{\partial p_{i_1\ldots i_\sigma}}\frac{\partial}{\partial q_{j_1\ldots j_{n-\tau}}}$에는 잘 알려져 있듯이 실제 의미 $\frac{\partial^2}{\partial p_{i_1\ldots i_\sigma}\partial q_{j_1\ldots j_{n-\tau}}}$가 대응한다. 열 $\frac{\partial}{\partial p_\sigma}$, $\frac{\partial}{\partial q_{n-\tau}}$는 열 $S^\sigma,T_{n-\tau}$에 동변하므로 이들과 같은 항등식들을 만족하며, 특히 (42.)와 (42a.), (43.)와 (43a.) 사이에 성립하는 항등식들도 만족한다. 요약하면 다음을 얻는다.

\emph{정리 VI:} 임의의 기본형들의 모든 불변 구성식은 이 기본형들에 기호적 수축 과정 (41.)을 반복 적용하거나, 다음 불변 미분 과정을 반복 적용하여 생긴다.
\begin{equation}
 \left(q_{n-\sigma-\lambda}\frac{\partial}{\partial p_\sigma}\middle|\frac{\partial}{\partial q_{n-\tau}}p_{\tau-\lambda}\right),
 \qquad(\sigma\geqq\tau,\ \lambda=1,2,\ldots,\tau,n-\sigma).
\tag{45}
\end{equation}
각 수축 (41.)에서 형식의 전체 무게, 즉 개별 변수열들의 무게의 합(\S\ 1 참조)은 양수 $2(\lambda^2+\lambda(\sigma-\tau))$만큼 감소한다는 점도 언급해야 한다. 따라서 일반적인 유한성 증명과 독립적으로 형식열, 즉 주어진 기본형의 계수에 관해 선형인 불변 구성식 전체의 유한성이 따른다.\footnote{Clebsch, a. a. O. \S\ 11 및 12를 참조하라. 기본 공변식 계의 유한성.}

% inherited: noether_paper04_ko_translation_001_20260804\targets\Noether_P04_Korean_T06_U32_UNCHECKED.tex
% sha256: 01165205A8F187C8538321CBA50EFED54C41211F051F39E7A274660D4321128B
더 나아가 정리 VI는 조건 (12.)를 만족하지 않는 형식들에 대해서도 성립한다. 실제로 (3.)의 각 인자 $\pair{S^\rho}{p_\rho}^{m_\rho}$를 $m_\rho$개의 선형 인자들의 곱 $\pair{A^\rho}{p_\rho}\pair{B^\rho}{p_\rho}\cdots\pair{M^\rho}{p_\rho}$로 바꿀 수 있다. 그러면 불변 구성식들은 선형형식 계의 불변 구성식들로 바뀌며, 이 선형형식들은 나중에 서로 같다고 놓기만 하면 된다. 그런데 각 개별 선형형식에 대해서는 조건 (12.)가 언제나 충족된다. 미분기호를 도입했을 때 이 조건들이 2계 미분계수만을 포함하기 때문이다.

% inherited: noether_paper04_ko_translation_001_20260804\targets\Noether_P04_Korean_T07_U33_UNCHECKED.tex
% sha256: A4FE355CA44A9EC87E1F86287D3CE18834FF8ADD7966DA00A45C59FFA5BEC8D5
\section*{\S\ 7. 일반 형식의 정규형에 의한 전개.}
이하(\S\ 8)에서 우리는 수축 과정 (41.)의 한 가지 주된 성질, 곧 그것이 기본 수축들로부터 합성될 수 있다는 사실을 유도할 것이다. 그러나 이를 위해서는 Mertens가\footnote{F. Mertens, \emph{\"Uber invariante Gebilde quatern\"arer Formen} (서론 참조), 이하 “Mertens”로 인용한다.} 4항 영역에서 제시한 중요한 정리를 $n$항 영역으로 옮길 필요가 있다. 이 정리에 따르면 임의의 유한한 형식계는 그와 동치인 정규형의 계로 바꿀 수 있다. 여기서 우리는 --- 이항 및 삼항 정의를 일반화하여\footnote{Study, a. a. O. S. 54.} --- “정규형”을 기본형 자체를 제외하고 계수에 관해 선형인 모든 불변 구성식이 항등적으로 소멸하는 형식으로 이해한다. 따라서 정리 VI(\S\ 6)에 의해 정규형 $\varphi$는 다음 미분방정식계를 만족한다.
\begin{equation}
 \left(q_{n-\sigma-\lambda}\frac{\partial}{\partial p_\sigma}\middle|\frac{\partial}{\partial q_{n-\tau}}p_{\tau-\lambda}\right)\varphi=0,
 \qquad(\sigma\geqq\tau,\ \lambda=1,2,\ldots,\tau,n-\sigma),
\tag{46}
\end{equation}
여기서 열 $q_{n-\sigma-\lambda},p_{\tau-\lambda}$는 임의의 양으로 보아야 한다. 이 미분방정식들은 4항 영역에서 --- $\sigma=\tau=2$일 때 (39.)를 고려하면 --- Mertens가 임의의 열 $p,q$를 도입하지 않고 제시한 10개의 미분방정식(a. a. O., 공식 28))과 완전히 일치한다. 이 방정식들과 정리 VI를 알면 그곳의 증명법을 거의 그대로 $n$개 변수의 경우로 옮길 수 있다. 추가해야 할 유일한 것은 분해 항등식 (30.)으로 구성된 형식들이 정규형임을 보이는 증명이다. 4항 영역에서는 이 증명이 더 이상의 변형 없이도 얻어진다. 그러므로 증명을 간략히 서술하는 것으로도 충분하다. 또한 단순화를 위하여 이하에서 유일하게 고려해야 하는 하나의 기본형 $f$와 그 형식열의 경우(\S\ 6 끝 참조)만을 다룬다. 임의로 많은 기본형들의 계수에 관해 임의 차수인 불변 구성식에 대해서도 증명은 같기 때문이다.

증명의 기본 착상은 변환된 계수를 도입하여, $n-1$개의 열 $p_\rho$를 가진 한 형식을 $x$에 동변인 $n$개의 열 $\xi$, 곧 변환량의 열들을 가진 형식들로 바꾸는 것이다. 이 형식들로부터 불변 미분 과정에 의해 찾는 정규형들이 직접 유도된다. 일반 형식을 정규형으로 전개하는 것은 $x$에 동변인 $\rho$개의 열 $\xi$를 가진 형식들($\rho\leqq n$)에 대한 급수 전개를 통하여 이루어진다. 불변 미분 과정은 원래의 열 $p_\rho$를 가진 형식들로 되돌아가게 한다.

% inherited: noether_paper04_ko_translation_001_20260804\targets\Noether_P04_Korean_T07_U34_UNCHECKED.tex
% sha256: 29B5581122A9C993878976CEFEBEFC068DA729F56946FD63C3DB6B66B3784DE0
$\xi^{(1)},\ldots,\xi^{(n)}$을 변환량의 열이라 하자. 그러면 다음이 성립한다.
\begin{equation}
\begin{aligned}
x_i&=\xi_i^{(1)}x'_1+\xi_i^{(2)}x'_2+\cdots+\xi_i^{(n)}x'_n,\\
p_{i_1\ldots i_\rho}&=\sum_j(\xi_{i_1}^{(j_1)}\cdots \xi_{i_\rho}^{(j_\rho)})p'_{j_1\ldots j_\rho}.
\end{aligned}
\tag{47}
\end{equation}
형식 $f,\varphi,\ldots$의 변환된 계수들을 $(f),(\varphi),\ldots$로 나타내고, 다음과 같이 놓자.
\[
 f\equiv\pair{S^1}{p_1}^{m_1}\pair{S^2}{p_2}^{m_2}\cdots\pair{S^{n-1}}{p_{n-1}}^{m_{n-1}}.
\]
그러면 다음을 얻는다.
\begin{equation}
\begin{aligned}
(f)&=\pair{S^1}{\xi^{(1)}}^{m_{11}}\cdots\pair{S^1}{\xi^{(n)}}^{m_{1n}}
\pair{S^2}{\xi^{(1)}\xi^{(2)}}^{m_{21}}\cdots
\pair{S^{n-1}}{\xi^{(2)}\cdots\xi^{(n)}}^{m_{n-1,n}}\\
&=\pair{s^{(1)}}{\xi^{(1)}}^{k_1}\pair{s^{(2)}}{\xi^{(2)}}^{k_2}\cdots\pair{s^{(n)}}{\xi^{(n)}}^{k_n},
\end{aligned}
\tag{48}
\end{equation}
여기서 $m_{11}+\cdots+m_{1n}=m_1$ 등이 성립한다. 다음 조건을 만족하는 각 계수 $(f)$로부터
\begin{equation}
 k_1\geqq k_2\geqq k_3\geqq\cdots\geqq k_n,
\tag{49}
\end{equation}
열 $\xi$를 정확히 모두 소진하는 다음 미분 과정을 거치면
\begin{equation}
\left(\frac{\partial}{\partial\xi^{(1)}}\middle|p_1\right)^{k_1-k_2}
\left(\frac{\partial}{\partial\xi^{(1)}}\frac{\partial}{\partial\xi^{(2)}}\middle|p_2\right)^{k_2-k_3}
\cdots
\left(\frac{\partial}{\partial\xi^{(1)}}\cdots\frac{\partial}{\partial\xi^{(n-1)}}\middle|p_{n-1}\right)^{k_{n-1}-k_n}
\left(\frac{\partial}{\partial\xi^{(1)}}\cdots\frac{\partial}{\partial\xi^{(n)}}\right)^{k_n}
\tag{50}
\end{equation}
다음 형식을 얻는다.
\begin{equation}
 \varphi=\bigl(s^{(1)}\mid p_1\bigr)^{k_1-k_2}
 \bigl(s^{(1)}s^{(2)}\mid p_2\bigr)^{k_2-k_3}\cdots
 \bigl(s^{(1)}\cdots s^{(n-1)}\mid p_{n-1}\bigr)^{k_{n-1}-k_n}
 \bigl(s^{(1)}s^{(2)}\cdots s^{(n)}\bigr)^{k_n}.
\tag{51}
\end{equation}
형식 $\varphi$들이 바로 찾던 정규형들이다. 실제로 이들은 정규형을 정의하는 다음 성질들을 가진다.

% inherited: noether_paper04_ko_translation_001_20260804\targets\Noether_P04_Korean_T07_U35_UNCHECKED.tex
% sha256: 2A5295BCFFD02A0F4874257C79FC289B37030989D839EFAF942CE6CDA3FABE1F
1. 이들은 기본형 $f$의 계수에 관해 선형인 불변 구성식이다. 이는 이들이 $f$에 불변 미분 과정 $\left(\frac{\partial}{\partial p_\rho}\middle|\xi^{(i_1)}\xi^{(i_2)}\cdots\xi^{(i_\rho)}\right)$ 및 (50.)을 적용하여 생기기 때문이다. 따라서 이들은(\S\ 6 끝에 따라) $f$의 유한 형식열의 편극들의 선형식으로 나타낼 수 있다. 이러한 편극들은 $\varphi$에 들어가는 변수 $p_\rho$를, 변수열 $p_\rho$를 가진 다음 형식, 곧 선형형식들의 곱으로 분해되는 형식의 계수로 간주하여 얻는다.
\begin{equation}
 H=(q_{n-1}\mid p_{n-1})^{k_1-k_2}(q_{n-2}\mid p_{n-2})^{k_2-k_3}\cdots(q_1\mid p_1)^{k_{n-1}-k_n}.
\tag{52}
\end{equation}
형식열 $f$와 $H$의 연립불변식들은 고려해야 하는 편극들을 모두 준다. 이들 편극은 각 개별 열 $p_\rho$에 관하여 $\varphi$ 자체와 같은 차수를 가져야 하기 때문이다.

2. 이들은 미분방정식 (46.)을 만족한다. 실제로 $\varphi$에 미분 과정 (45.)을 적용하면 다음이 나온다.
\[
 \varphi'=\bigl(q_{n-\sigma-\lambda}s^{(1)}\cdots s^{(\sigma)}\mid [s^{(1)}\cdots s^{(\tau)}]_{n-\tau}p_{\tau-\lambda}\bigr),\qquad (\sigma\geqq\tau).
\]
분해 항등식 (30.)에서 열 $s$로 만들어진 열 $q_\sigma=s^{(1)}\cdots s^{(\sigma)}$, $p_{n-\tau}=[s^{(1)}\cdots s^{(\tau)}]_{n-\tau}$를 그곳의 $q_\rho,p_\tau$와 동일시하면 다음을 얻는다.
\[
 \varphi'=\sum_{\beta=\sigma-\tau+\lambda}^{n-\tau,\sigma}\sum_{i_\beta}\eps\,
 \pair{\dot{q}_{\sigma-\beta}^{(i_1\ldots i_\beta)}q_{n-\tau+\lambda}}{p_{\sigma+\lambda}\dot{p}_{n-\tau-\beta}^{(i_1\ldots i_\beta)}}.
\]
이때 다음이 성립한다.
\[
 \dot{p}_{n-\tau-\beta}^{(i_1\ldots i_\beta)}\sim s^{(1)}\cdots s^{(\tau)}s^{(i_1)}\cdots s^{(i_\beta)},
\]
여기서 $i_1,\ldots,i_\beta$는 1부터 $\sigma$까지의 수 가운데서 고른, 서로 다른 임의의 $\beta$개 수를 뜻한다. 이제 $\beta>\sigma-\tau$이므로 이 수들 $i_1,\ldots,i_\beta$ 중에는 반드시 1부터 $\tau$까지의 수가 들어 있다. 그러면 $\dot{p}_{n-\tau-\beta}^{(i_1\ldots i_\beta)}$는 항등적으로 소멸하고, 따라서 각 개별 행렬곱 및 $\varphi'$도 소멸한다.

% inherited: noether_paper04_ko_translation_001_20260804\targets\Noether_P04_Korean_T07_U36_UNCHECKED.tex
% sha256: 1E7EEDA0FF9B95F51C6DDE6E22A5E37257EC4A175DA6EB081973D263BE035020
정규형의 계가 형식열 $f$와 동치임을 보이기 위해 이제 남은 것은, 형식열의 모든 형식을 1.에서 정한 의미에서 정규형들의 편극들에 따라 전개할 수 있음을 보이는 것뿐이다. $\Pi$를 다음 편극 연산들의 결합식이라 하자.
\srcnumdisplay[4.2em]{(53.)}{%
 \left(\frac{\partial}{\partial\xi^{(i)}}\middle|\xi^{(k)}\right),
 \qquad (i\ne k;\ i=1,2,\ldots,\rho;\ k=1,2,\ldots,\rho)}
또한 $\Pi^\sigma$를 $i=\sigma$이고 $k<\sigma$인 특수 연산 (53.)들의 결합식이라 하자. 열 $\xi^{(\rho)}$에 관해 차수가 $k_\rho$인, $\rho$개의 열 $\xi^{(1)},\ldots,\xi^{(\rho)}$의 형식 $F$에 대해서는 다음 전개가 성립한다\srcfn{*)}{F. Mertens, \emph{\"Uber eine Formel der Determinantentheorie}. Sitzungsber. d. Ak. d. Wiss., Wien, math.-naturw. Kl., Bd. 91. 2. Abt. (1885). Formel 20).}:
\srcnumdisplay[4.2em]{(54.)}{%
 F=\sum_{\alpha=0}^{k_\rho}\Pi F_\alpha,
 \qquad(\rho\leqq n)}
여기서 $F_\alpha$는 마지막 열 $\xi^{(\rho)}$에 관해 $\alpha$차이고, $F$에 다음 불변 미분 연산
\srcnumdisplay[4.2em]{(55.)}{%
 \left(\frac{\partial}{\partial\xi^{(1)}}\cdots\frac{\partial}{\partial\xi^{(\rho)}}\middle|\xi^{(1)}\cdots\xi^{(\rho)}\right)^\alpha}
및 $\Pi^\rho$를 적용하여 생긴다. $F_\alpha$를 생성할 때 과정 (55.)를
\srcnumdisplay[4.2em]{(56.)}{%
 \left(\frac{\partial}{\partial\xi^{(1)}}\cdots\frac{\partial}{\partial\xi^{(\rho)}}\middle|p_\rho\right)^\alpha,}
로 바꾸면, $\rho-1$개의 열 $\xi^{(1)},\ldots,\xi^{(\rho-1)}$만을 가진 형식 $F_\alpha(p_\rho)$가 생기며, 여기에 다시 (54.)를 적용할 수 있다. 이 절차를 계속하면 형식 $G(p_\rho,p_{\rho-1},\ldots,p_1)$에 이른다. 이 형식들로부터 (54.)에 따른 $F$의 전개를 얻으려면, 각 열 $p_\sigma$를 $\xi^{(1)},\ldots,\xi^{(\sigma)}$로 바꾸고 생기는 형식들에 편극 과정 $\Pi$ (53.)를 적용해야 한다. 그러면 (48.)과 비교할 때 변환된 계수 $(G)$들의 합을 얻는다. $\rho=n$이면 첫 단계에서 과정 (55.)를 그대로 둔다. $c$가 수치 상수들을 뜻한다고 하고, 더 나아가 다음과 같이 줄여 쓰자.
\srcnumdisplay[4.2em]{(57.)}{%
 (\xi^{(1)}\xi^{(2)}\cdots\xi^{(n)})=\Delta;
 \qquad
 \left(\frac{\partial}{\partial\xi^{(1)}}\frac{\partial}{\partial\xi^{(2)}}\cdots\frac{\partial}{\partial\xi^{(n)}}\right)=\nabla,}
그러면 $\rho=n$인 경우에는
\srcnumdisplay[4.2em]{(58.)}{%
 F=\sum c\Delta^\alpha(G),}
을 얻는 반면, $\rho<n$이면 오른쪽에는 $\alpha=0$인 항, 곧 $\sum c(G)$만 나타난다.

% inherited: noether_paper04_ko_translation_001_20260804\targets\Noether_P04_Korean_T07_U37_UNCHECKED.tex
% sha256: 8F2DC0D514194A08AA3E0999B0C08ECBD115A3FB8259558BB65045AEF186C527
(58.)을 변환된 계수 $(f)$ (48.)에 적용하자. 편극 과정 $\Pi^\sigma$는 $\rho\geqq\sigma$인 모든 연산 (56.)과 교환할 수 있고, 변환된 계수의 편극은 다시 변환된 계수가 되므로 다음을 얻는다.
\[
\begin{aligned}
 G={}&\left(\frac{\partial}{\partial\xi^{(1)}}\middle|p_1\right)^{\beta_1}
 \left(\frac{\partial}{\partial\xi^{(1)}}\frac{\partial}{\partial\xi^{(2)}}\middle|p_2\right)^{\beta_2}
 \cdots
 \left(\frac{\partial}{\partial\xi^{(1)}}\cdots\frac{\partial}{\partial\xi^{(n-1)}}\middle|p_{n-1}\right)^{\beta_{n-1}}
 \nabla^{\beta_n}\Pi(f)\\
 ={}&\sum c\varphi
\end{aligned}
\]
따라서 (58.)로부터 다음이 따른다.
\srcnumdisplay[4.2em]{(59.)}{%
 (f)=\sum c\Delta^\alpha(\varphi)\qquad\text{(Mertens, Formel 23)).}}
(59.)에서 $f$ 자체의 전개로 넘어가기 위해, 다시 변수 $p_\rho$를 $f$에 대응하는 차수 $m_1,\ldots,m_{n-1}$을 가진 형식 $H$ (52.)의 계수로 간주하고, $m=m_1+\cdots+m_{n-1}$이라 하자. 그러면 불변식의 정의에 따라 다음이 성립한다.
\[
 f\cdot\Delta^m=\sum(f)(H)=\sum c\Delta^\alpha(\varphi)(H),
\]
열 $\xi$를 정확히 모두 소진하는 과정 $\nabla^m$을 적용하면 다음을 얻는다.
\srcnumdisplay[4.2em]{(60.)}{%
 f=\sum c\Phi,}
여기서 형식 $\Phi$는 $\varphi$와 $H$로부터 변수들에 관한 불변 미분 과정을 통해 유도되므로, 형식열 $\varphi$와 $H$의 연립불변식이다\srcfn{*)}{직접 이 결론을 내리는 대신 Mertens, \S\ 7에서는 본질적으로 형식열 $H$를 만드는 데 목적을 둔 일련의 추가 미분 과정들을 도입하는데, 이는 우리의 정리 VI가 수행하는 일이다.}. 그런데 $\varphi$가 정규형이므로 형식열 $\varphi$는 형식 $\varphi$ 자체로 제한되고, 형식열 $H$는 가능한 모든 불변 변수 결합을 포함한다. 따라서 $\varphi$와 $H$의 연립불변식들은 1.에서 말한 의미의 $\varphi$의 편극들이다. 즉, (60.)은 $f$를 정규형들의 편극들에 따라 전개한 것이다. 이제 형식열 $f$의 임의의 형식 $\theta$에 대해서도 같은 전개를 증명해야 한다. $\Gamma$를 $(\theta)$로부터 (50.)에 의해 유도된 정규형들이라 하자. 그러면 다음이 성립한다.

% inherited: noether_paper04_ko_translation_001_20260804\targets\Noether_P04_Korean_T07_U38_UNCHECKED.tex
% sha256: DDF04812D84E68452CBBCB0163A4C501E64091FEAB369C2FD9B7D21510B80B94
\srcnumdisplay[4.2em]{(61.)}{%
 (\theta)=\sum c\Delta^\beta(\Gamma).}
형식 $\theta$는 각 변환된 계수에 대해 성립하는 다음 관계로 특징지어진다.
\[
 (\theta)\Delta^\sigma=\sum c(f)\qquad\text{(Mertens, Formel 9)).}
\]
그런데 $\Delta^\sigma$를 인자로 포함하는 $n$개의 열 $\xi$의 임의의 형식 $F$는 둘째 것인 $\nabla^\sigma\Pi F$도 포함하므로(Mertens, Formel 20)), 다음이 따른다.
\[
 (\theta)=\sum c\nabla^\sigma(f),\qquad\text{따라서}\qquad \Gamma=\sum c\varphi,
\]
그리고 (61.)로부터 (59.)에 대응하는 다음 관계를 얻는다.
\[
 (\theta)=\sum c\Delta^\beta(\varphi),
\]
이 관계는 $\theta$에 대해 정규형 $\varphi$의 편극들에 따른 전개 (60.)을 이끌어 낸다. 요약하면 다음을 얻는다.

\emph{정리 VII. 모든 형식열 $f$는 그와 동치인 정규형의 계로 바꿀 수 있다.}

% inherited: noether_paper04_ko_translation_001_20260804\targets\Noether_P04_Korean_T08_U39_UNCHECKED.tex
% sha256: D5F068657C9E22A823C55E865C63990E559102F90329F9F1717A8313B03F6BEE
\begin{center}
\emph{\S\ 8. 수축 과정의 기본 수축으로의 환원.}
\end{center}
정리 VII에 이어 이제 \S\ 7의 첫머리에서 언급한 수축 과정의 기본 수축으로의 환원을 수행한다. 이때(\S\ 5 참조) 서로 동변인 두 열 $S^\rho,T^\rho$의 모든 수축을 “\emph{동변 수축}”이라 부르고, 특히 두 열 $S^\rho,T^\rho$의 결손 1인 동변 수축을 “\emph{유형 $\rho$의 기본 수축}”이라 부른다. 그러면 정리 VI를 보완하여 다음을 증명한다.

\emph{정리 VIII: 일반 수축 과정 (41.)은 $\rho=1,2,\ldots,n-1$인 유형 $\rho$의 $(n-1)$개 기본 수축으로 합성할 수 있다. 다시 말해, 모든 불변 구성식을 생성하려면 $n-1$개의 기본 수축을 차례로 적용하는 것으로 충분하다.}

이항 영역에서는 수축 과정 (41.)이 가능한 유일한 기본 수축과 곧바로 일치한다. 삼항 영역에서는 내가 학위논문의 \S\ 1에서(서론 참조) 정리 VIII을 증명하였다. 그곳에서는 (44.) 때문에 일반적인 동변 수축들이 여전히 기본 수축들과 동일하다. 주목할 점은 정리 VIII이 $n$개 변수에 대하여 기본 수축들로의 환원을 이루어 내며, 훨씬 약한 동변 수축들로의 환원에 그치지 않는다는 것이다. 증명은 원리상 삼항 경우를 본뜬 것으로, 기호적 항등식들을 적용하여 가장 일반적인 수축들을 기본 수축들로부터 구성한다. 우리는 다음을 생성할 것이다.
\begin{enumerate}
\renewcommand{\labelenumi}{\arabic{enumi})}
\item 결손 1의 임의의 수축을 결손 1의 동변 수축, 즉 기본 수축으로부터 생성한다(삼항 경우의 대응물),
\item 결손 $\lambda$의 임의의 수축을 결손 $\lambda$의 동변 수축과 결손 1의 임의의 수축으로부터 생성한다,
\item 결손 $\lambda$의 동변 수축을 결손 $\lambda-1$의 동변 수축과 결손 1의 임의의 수축으로부터 생성한다.
\end{enumerate}
이때 본질적으로 중요한 것은 정리 VII에 따라 수축하려는 두 형식 $S$와 $T$를 정규형으로 가정한다는 점이다. 따라서 (46.)에 따라 다음 관계들이 성립한다.
\srcnumdisplay[4.2em]{\raisebox{-1.7ex}{(62.)}}{%
\begin{aligned}
 \pair{q_{n-\sigma_1-\lambda}S^{\sigma_1}}{S_{n-\sigma_2}p_{\sigma_2-\lambda}}&=0,
 && (\sigma_1\geqq\sigma_2,\ \lambda=1,2,\ldots,\sigma_2,n-\sigma_1),\\
 \pair{q_{n-\tau_1-\lambda}T^{\tau_1}}{T_{n-\tau_2}p_{\tau_2-\lambda}}&=0,
 && (\tau_1\geqq\tau_2,\ \lambda=1,2,\ldots,\tau_2,n-\tau_1).
\end{aligned}}
더 나아가 \S\ 2의 끝부분에 따라 형식 $S$와 $T$가 모든 변수열에 관해 선형이라고 가정하면 충분하다. 즉, 구성되는 형식들은 다음 형식열에 속한다.
\[
 f=\pair{S^1}{p_1}\pair{S^2}{p_2}\cdots\pair{S^{n-1}}{p_{n-1}}
   \pair{T^1}{p_1}\pair{T^2}{p_2}\cdots\pair{T^{n-1}}{p_{n-1}}.
\]

% inherited: noether_paper04_ko_translation_001_20260804\targets\Noether_P04_Korean_T08_U40_UNCHECKED.tex
% sha256: 9CD0BF3111D5239DA9B4E60F8F8B1E877F551A962485628A9E23866A80718A98
\emph{1) $\sigma>\tau$일 때 $\pair{q_{n-\sigma-1}S^\sigma}{T_{n-\tau}p_{\tau-1}}$를 유형 $\tau+\alpha$의 기본 수축들로부터 생성하기; $\alpha=0,1,\ldots,\sigma-\tau$.} 유형 $\tau$의 기본 수축
\[
 \pair{q_{n-\tau-1}S^\tau}{T_{n-\tau}p_{\tau-1}}
\]
에서 대입
\[
 w\sim T_{n-\tau}p_{\tau-1}
\]
을 통해 위 무게지수가 $\tau+1$인 세 열을 가진 형식열 $f$의 한 형식을 얻는다(\S\ 1 참조).
\srcnumdisplay[4.2em]{(63.)}{%
 \pair{wS^\tau}{p_{\tau+1}}\pair{S^{\tau+1}}{p_{\tau+1}}\pair{T^{\tau+1}}{p_{\tau+1}}.}
(63.)에 유형 $\tau+1$의 기본 수축을 적용하면 다음을 얻는다.
\[
 \pair{q_{n-\tau-2}wS^\tau}{S_{n-\tau-1}p_\tau};
\]
여기에서 (31.)과 대입 $q_{n-\tau-2}w=q'_{n-\tau-1}$을 통해 다음 값을 얻는다.
\[
 \eps\pair{q_{n-\tau-2}wS^{\tau+1}}{S^\tau p_\tau}
 +\sum_{\alpha=1}^{\tau,n-\tau-1}\Pi\pair{q'_{n-\tau-1-\alpha}S^{\tau+1}}{S_{n-\tau}p_{\tau-\alpha}}.
\]
합 기호 아래의 식은 (62.)에 따라 항등적으로 소멸한다. 따라서 다음 형식에 이른다.
\[
 \pair{wS^{\tau+1}}{p_{\tau+2}}\pair{S^{\tau+2}}{p_{\tau+2}}\pair{T^{\tau+2}}{p_{\tau+2}},
\]
이는 (63.)에서 $\tau$를 $\tau+1$로 바꾸어 생기는 형식이다. 그러므로 이 과정을 $(\sigma-\tau)$번 반복하면 찾던 수축에 이른다.
\[
 \pair{wS^\sigma}{p_{\sigma+1}}=\eps\pair{q_{n-\sigma-1}S^\sigma}{T_{n-\tau}p_{\tau-1}}.
\]
방금 수행한 과정을 다음 공식으로 간단히 나타낸다.
\srcnumdisplay[4.2em]{(64.)}{%
 T^\tau S^\tau,\quad S^{\tau+1},\quad S^{\tau+2},\ldots,S^\sigma;}
여기서는 $S$의 정규형 성질만 사용되었고 $T$의 정규형 성질은 사용되지 않았음을 알 수 있다. (64.)와 쌍대인 과정은 $T$의 정규형 성질만 사용하여 역순으로 수행할 수 있다. 곧 다음 과정이다.
\srcnumdisplay[4.2em]{(65.)}{%
 S^\sigma T^\sigma,\quad T^{\sigma-1},\quad T^{\sigma-2},\ldots,T^\tau}
이는 다음 관계들에 따른다.
\begin{align*}
\pair{q_{n-\sigma-1}S^\sigma}{T_{n-\sigma}p_{\sigma-1}}&=\eps\bigl(T_{n-\sigma}y p_{\sigma-1}\bigr),\quad \text{여기서 }y\sim q_{n-\sigma-1}S^\sigma,\\
\pair{q_{n-\sigma}T^{\sigma-1}}{T_{n-\sigma}yp_{\sigma-2}}&=\eps\bigl(T_{n-\sigma+1}y p_{\sigma-2}\bigr),\\
&\vdots\\
\pair{q_{n-\tau-1}T^\tau}{T_{n-\tau-1}yp_{\tau-1}}&=\eps\pair{q_{n-\sigma-1}S^\sigma}{T_{n-\tau}p_{\tau-1}}.
\end{align*}
끝으로 \S\ 6의 끝부분에서 제시한 전체 무게의 감소와의 관련성을 지적해 두자. 그 값은 결손 1의 수축에서는 $2(1+(\sigma-\tau))$이고, 기본 수축에서는 $2\cdot1$이다. 따라서 기본 수축들로의 합성이 가능하기만 하다면, 위에서 수행한 것처럼 반드시 $\sigma-\tau+1$개의 기본 수축이 필요하다.

% inherited: noether_paper04_ko_translation_001_20260804\targets\Noether_P04_Korean_T08_U41_UNCHECKED.tex
% sha256: C943587E0D0AB85E9DE341529F005C12DA8E6AED8B123F222C3B62FFB5FD316E
\emph{2) 일반 수축을 동변 수축과 기본 수축들로부터 생성하기.} 일반 수축 (41.)의 무게 감소는 $2(\lambda^2+\lambda(\sigma-\tau))=2[\lambda^2+(\sigma-\tau)(1+(\lambda-1))]$이다. 따라서 이를 결손 $\lambda$의 동변 수축(무게 감소 $2\lambda^2$)과, 결손이 1이고 무게지수의 차가 $\lambda-1$인 $\sigma-\tau$개의 수축(무게 감소 $2\{1+(\lambda-1)\}$)으로 분해할 가능성이 생긴다. $\lambda=1$인 경우 이 분해는 1)에서 제시한 분해와 일치한다. 이제 이 분해를 실제로 실현할 수 있음을 보이겠다.

결손 $\lambda$의 동변 수축
\[
 \pair{q_{n-\tau-\lambda}S^\tau}{T_{n-\tau}p_{\tau-\lambda}}
\]
에서 대입 $R^\lambda\sim T_{n-\tau}p_{\tau-\lambda}$을 통해 위 무게지수가 각각 $\tau+\lambda$와 $\tau+1$인 두 열을 가진 형식을 얻는다.
\srcnumdisplay[4.2em]{(66.)}{%
 \pair{R^\lambda S^\tau}{p_{\tau+\lambda}}\pair{S^{\tau+1}}{p_{\tau+1}}.}
더 낮은 무게지수 $\tau+1$을 가진 열은 정규형에 속한다. 따라서 (66.)에는 (65.)의 순서에 따라 $\lambda$개의 기본 수축으로 합성된 결손 1의 수축을 적용할 수 있다.
\[
 (R^\lambda S^\tau)S^{\tau+\lambda},\quad S^{\tau+\lambda-1},\quad S^{\tau+\lambda-2},\ldots,S^{\tau+1}.
\]
(31.)과 (62.)를 고려하면 이는 다음을 준다.
\[
 \pair{q_{n-\tau-\lambda-1}R^\lambda S^\tau}{S_{n-\tau-1}p_\tau}=\eps\pair{R^\lambda S^{\tau+1}}{p_{\tau+\lambda+1}},
\]
즉, 1)에서 (63.)에 대해 그러했던 것과 마찬가지로 (66.)에서 $\tau$를 $\tau+1$로 바꾸어 생기는 형식을 준다. 이 과정을 $(\sigma-\tau)$번 반복하면 찾던 수축에 이른다.
\[
 \pair{R^\lambda S^\sigma}{p_{\sigma+\lambda}}=\eps\pair{q_{n-\sigma-\lambda}S^\sigma}{T_{n-\tau}p_{\tau-\lambda}}.
\]
이 전체 과정에서는 $S$의 정규형 성질만 사용되었다. $T$를 정규형으로 보면 순서가 완전히 뒤집힌 쌍대 과정에 이르며, 그 공식들은 다음과 같다.
\begin{align*}
\pair{q_{n-\sigma-\lambda}S^\sigma}{T_{n-\sigma}p_{\sigma-\lambda}}&=\eps\bigl(T_{n-\sigma}R_\lambda p_{\sigma-\lambda}\bigr),\quad \text{여기서 }R_\lambda\sim q_{n-\sigma-\lambda}S^\sigma,\\
\pair{q_{n-\sigma}T^{\sigma-1}}{T_{n-\sigma}R_\lambda p_{\sigma-\lambda-1}}&=\eps\bigl(T_{n-\sigma+1}R_\lambda p_{\sigma-\lambda-1}\bigr),\\
&\vdots\\
\pair{q_{n-\tau-1}T^\tau}{T_{n-\tau-1}R_\lambda p_{\tau-\lambda}}&=\eps\pair{q_{n-\sigma-\lambda}S^\sigma}{T_{n-\tau}p_{\tau-\lambda}}.
\end{align*}

% inherited: noether_paper04_ko_translation_001_20260804\targets\Noether_P04_Korean_T08_U42_UNCHECKED.tex
% sha256: FC0E0CEF4FE64A66E0D928A11433602B51230FC91B88F6797681955F5E99BF5E
\emph{3) 결손 $\lambda-1$의 동변 수축과 기본 수축들로부터 결손 $\lambda$의 동변 수축을 생성하기.} 결손 $\lambda$의 동변 수축에서 무게 감소는 $2\lambda^2=2((\lambda-1)^2+2\lambda-1)$이므로, 이를 결손 $\lambda-1$의 동변 수축과 $2\lambda-1$개의 기본 수축으로 분해할 가능성이 생긴다. 먼저 특별히 $n=2\lambda$로 놓으면 이 분해를 가장 쉽게 찾을 수 있다. 유일하게 가능한 결손 $\lambda$의 동변 수축은 $\pair{S^{n/2}}{T_{n/2}}=\pair{S^\lambda}{T_\lambda}$가 된다. 이 수축은 같은 열들의 결손 $\lambda-1$인 수축 $\pair{uS^\lambda}{T_\lambda x}$보다 열 $u$와 $x$를 하나씩 덜 포함한다. 역으로 결손 $\lambda$의 수축을, 열 $u$와 $x$를 소멸시키고(무게 감소 $2(n-1)=2(2\lambda-1)$) 동시에 새로운 기호열 $R^\lambda$를 생성하는 기본 수축들과 결손 $\lambda-1$의 동변 수축으로 합성한다. 이를 수행하는 가장 단순한 수축은 다음 결손 1의 수축이다.
\[
 \pair{sT^\lambda}{\sigma p_\lambda}=\eps\pair{S^{2\lambda-1}}{R_{\lambda-1}p_\lambda}=\eps\pair{R^\lambda}{p_\lambda},
\]
여기서 다음과 같이 놓았다.
\[
 R^{\lambda+1}=sT^\lambda,
 \quad s=S^1,
 \quad \sigma=S_1,
 \quad R^\lambda\sim\sigma R_{\lambda-1}.
\]
이 수축은 (65.)와 (64.)에 따라 다음 $2\lambda-1$개의 기본 수축으로 합성된다.
\srcnumdisplay[4.2em]{(67.)}{%
 T^\lambda S^\lambda,S^{\lambda-1},\ldots,S^1;
 \qquad R^{\lambda+1}S^{\lambda+1},S^{\lambda+2},\ldots,S^{2\lambda-1}.}
결손 $\lambda-1$의 수축을 적용하면 (21.)과 (62.)를 고려하여 다음을 얻는다.
\[
 \pair{uS^\lambda}{\sigma R_{\lambda-1}x}=\eps\pair{R^{\lambda+1}}{S_\lambda x}
 =\eps\pair{sT^\lambda}{S^\lambda x}=\eps\pair{S^\lambda}{T_\lambda}.
\]

% inherited: noether_paper04_ko_translation_001_20260804\targets\Noether_P04_Korean_T08_U43_UNCHECKED.tex
% sha256: A0E31D058D3FD8B86F05CED6278E14124714B1019CBAC6ACB99614DD45616CE6
이제 두 열 $S^\rho,T^\rho$와 임의의 $n$에 관한 일반적인 경우는 특수한 경우로부터 직접 옮길 수 있다. (67.) 대신 다음 $2\lambda-1$개의 기본 수축이 들어간다.
\srcnumdisplay[4.2em]{(68.)}{%
 T^\rho S^\rho,S^{\rho-1},\ldots,S^{\rho-\lambda+1};
 \qquad R^{\rho+1}S^{\rho+1},S^{\rho+2},\ldots,S^{\rho+\lambda-1}.}
수축들의 전반부는 다음 형식을 준다.
\[
 \pair{q_{n-\rho-1}T^\rho}{S_{n-\rho+\lambda-1}p_{\rho-\lambda}},
\]
여기에 대입
\[
 S_{n-\rho+\lambda-1}p_{\rho-\lambda}\sim w,
 \qquad T^\rho w=R^{\rho+1}
\]
을 하고 수축들의 후반부를 적용하면 다음 형식이 생긴다.
\[
 \pair{q_{n-\rho-\lambda}S^{\rho+\lambda-1}}{R_{n-\rho-1}p_\rho}.
\]
다음과 같이 대입하면
\[
 q_{n-\rho-\lambda}S^{\rho+\lambda-1}\sim y,
\]
결손 $\lambda-1$의 동변 수축에 의해 다음을 얻는다.
\srcnumdisplay[4.2em]{(69.)}{%
 \pair{q_{n-\rho-\lambda+1}S^\rho}{yR_{n-\rho-1}p_{\rho-\lambda+1}}.}
이것이 찾던 결손 $\lambda$의 수축임을 보이자. 다음과 같이 대입하고
\[
 R_{n-\rho-1}p_{\rho-\lambda+1}=R_{n-\lambda}
\]
$y$를 풀어 놓으면 다음을 얻는다는 점을 고려하면
\[
 \pair{q_{n-\rho-\lambda+1}R^\lambda}{S_{n-\rho}y}
 =\eps\pair{q_{n-\rho-\lambda}S^{\rho+\lambda-1}}{S_{n-\rho}p'_{\rho-1}}=0
\]
(20.)에 따라 (69.)의 값은 다음과 같다.
\[
 \eps\pair{S^\rho R^\lambda}{p_{\rho+\lambda-1}y}
 =\eps\pair{q_{n-\rho-\lambda}S^{\rho+\lambda-1}}{[S^\rho R^\lambda]_{n-\rho-\lambda}p_{\rho+\lambda-1}}.
\]
그런데 이 수축은 유형 (43.)이므로 다음과 동치이다.
\[
 \pair{q_{n-\rho-\lambda}S^\rho}{R_{n-\rho-1}p_{\rho-\lambda+1}}
 =\eps\pair{wT^\rho}{R_\lambda p_{\rho-\lambda+1}},
 \qquad \text{여기서 }R_\lambda\sim q_{n-\rho-\lambda}S^\rho.
\]
여기서 $R_\lambda$는 이미 원하는 종결형을 갖고 있으며, 이 식은 특수한 경우의 형식 $\pair{sT^\lambda}{S_\lambda x}$에 대응한다. 이제 $w$와 $R_\lambda$를 풀어 놓으면 다음을 얻는다는 점을 고려하자.
\[
 \pair{wR^{n-\lambda}}{T_{n-\rho}p_{\rho-\lambda+1}}
 =\eps\pair{q'_{\lambda-1}R^{n-\lambda}}{S_{n-\rho+\lambda-1}p_{\rho-\lambda}}
 =\eps\pair{q''_{n-\rho-1}S^\rho}{S_{n-\rho+\lambda-1}p_{\rho-\lambda}}=0,
\]
그러면 (21.)에 따라 다음 값을 얻는다.
\[
 \pair{wq_{n-\rho+\lambda-1}}{T_{n-\rho}R_\lambda}
 =\eps\pair{q_{n-\rho+\lambda-1}[T_{n-\rho}R_\lambda]^{n-\lambda}}{S_{n-\rho+\lambda-1}p_{\rho-\lambda}}
 \quad\text{동치}\quad
 \pair{q_{n-\rho-\lambda}S^\rho}{T_{n-\rho}p_{\rho-\lambda}}.
\]
이 전체 과정에서는 $S$의 정규형 성질만 사용되었고 $T$의 정규형 성질은 사용되지 않았다. 따라서 (69.)를 생성하기 위한 결손 $\lambda-1$의 동변 수축은 $2\lambda-3$개의 기본 수축과 결손 $\lambda-2$의 동변 수축으로 바꿀 수 있고, 이 후자의 수축도 다시 기본 수축들로 분해할 수 있으며, 이와 같이 계속할 수 있다. 마찬가지로 $T$의 정규형 성질을 사용하고 쌍대적 과정, 즉 다음 기본 수축들을 적용하면 생성이 이루어진다.
\[
 S^\rho T^\rho,\quad T^{\rho+1},\ldots,T^{\rho+\lambda-1};
 \qquad R^{\rho-1}T^{\rho-1},\quad T^{\rho-2},\ldots,T^{\rho-\lambda+1},
\]
그리고 (69.)와 쌍대적인 수축을 적용한다. 그러면 앞의 수축과 동치인 다음 수축이 생긴다.
\[
 \pair{q_{n-\rho-\lambda}T^\rho}{S_{n-\rho}p_{\rho-\lambda}}.
\]

% inherited: noether_paper04_ko_translation_001_20260804\targets\Noether_P04_Korean_T08_U44_UNCHECKED.tex
% sha256: 9463F17C69AF95C2182C19281714DAA7C0DFFEF0F25CCC303C783D06504982CA
2)에서 얻은 분해를 덧붙이면 실제로 \emph{가장 일반적인 수축들을 기본 수축들로부터 생성하는 방법}을 얻는다. 이제 이 생성 방법이 순서의 교환을 제외하면 \emph{유일함}을 증명해야 한다. 즉, 각 수축에는 $\rho=1,2,\ldots,n-1$인 유형 $\rho$의 기본 수축 수 $\alpha_\rho$가 오직 하나만 대응함을 보여야 한다.

$m_\rho$를 출발 형식 $f$의 변수열 $p_\rho$에 관한 차수, $l_\rho$를 수축으로 생긴 형식의 차수라 하자. 유형 $\rho$의 각 기본 수축은 차수 $m_{\rho-1},m_\rho,m_{\rho+1}$을 각각 $m_{\rho-1}+1,m_\rho-2,m_{\rho+1}+1$로 바꾼다. 따라서 $\alpha$들은 다음 방정식계를 만족한다.
\srcnumdisplay[4.2em]{(70.)}{%
 m_\rho-l_\rho=-\alpha_{\rho-1}+2\alpha_\rho-\alpha_{\rho+1},
 \qquad (\rho=1,2,\ldots,n-1)}
여기서 $\alpha_0$와 $\alpha_n$은 0으로 놓아야 하며, 이 방정식계의 행렬식은 값 $n$을 갖는다\srcfn{*)}{$\Delta_n$이 $n$행 행렬식을 나타내면 다음 점화식이 성립한다: $\Delta_n=2\Delta_{n-1}-\Delta_{n-2}$; 즉 $\Delta_n-\Delta_{n-1}=\Delta_{n-1}-\Delta_{n-2}=\cdots=\Delta_2-\Delta_1=1$이다. 왜냐하면 $\Delta_2=3$, $\Delta_1=2$이기 때문이다.}. 따라서 $\alpha$들은 유일하게 결정되며, 정리 VIII에 따라 양의 정수이다. 이로부터 차 $m_\rho-l_\rho$에 대한 조건들이 생긴다.

이 절의 결과들은 \emph{관계 (62.)에 의거하여} 성립한다. 그러나 관계 (62.)가 성립하지 않는 경우에는 종결식들이 시작식들과 동치인 것이 아니다. 실제로 \S\ 6에서 제시한 동치의 정의는 다음과 같이도 표현할 수 있다\srcfn{**)}{이에 대해서는 다음 절 끝의 고찰을 참조하라.}: “두 형식이 더 높은 수축의 형식들, 즉 더 큰 무게 감소를 보이는 형식들만큼 서로 다를 때, 그리고 오직 그때에만 두 형식은 동치이다.” 이 더 큰 무게 감소는 수축에 들어가는 두 열 $q_{n-\sigma-\lambda},p_{\tau-\lambda}$가 실제로 변수열일 때에는 언제나 일어난다. 이는 (41.), (42.), 이 절에서 동치인 것으로 나타나는 수축들, 그리고 정리 VII의 동치인 형식들의 경우에도 그러했다. 그러나 이 열들이 기호열들로부터 대입을 통해 유도되었을 때에는 그러한 무게 감소가 일어날 필요가 없다. 관계 (62.)에 의거하여 서로 표현할 수 있었던 이 절의 형식들이 바로 그러한 경우였다. 소멸하는 형식들 가운데 일부는 오히려 더 큰 전체 무게를 가짐을 쉽게 알 수 있다.

% inherited: noether_paper04_ko_translation_001_20260804\targets\Noether_P04_Korean_T09_U45_UNCHECKED.tex
% sha256: A3855C2333A2C49CCC450510237915A722E6C1325DE0EB1B590B823183449311
\begin{center}
\emph{\S\ 9. 형식열. 환원 정리}\srcfn{*)}{삼항 영역에서의 대응하는 고찰을 참조하라. 학위논문 \S\S\ 1--3.}.
\end{center}
\S\ 6(끝)에서 도입한 \emph{형식열}을 이제 조금 더 엄밀하게 다음과 같이 정의한다. 즉 “주어진 출발형식의 계수에 관해 선형인 불변 구성식 전체, 다시 말해 출발형식을 자기 자신에 대해 수축하여 생긴 형식 전체를 \emph{더 높은 형식에 따른 순서로 배열한 것}”이다.

% inherited: noether_paper04_ko_translation_001_20260804\targets\Noether_P04_Korean_T09_U46_UNCHECKED.tex
% sha256: 46E36FAA42AFB76E8862B663F5F6C8FC95C37FB5EB37F7DD2CB02588E7C8D6EE
더 높은 형식을 설명하기 위해, 정리 VII에 따라 출발형식이 두 정규형의 곱 $S\cdot T$로 주어졌다고 하자. 그러면 자기 자신에 대한 더 높은 차수의 수축으로 생긴 형식을 \emph{더 높은 형식}이라 부르고, 자기수축의 차수가 같은 형식들 가운데서는 특별히 정규화한 형식을 그렇게 부른다. 더 정확히 말하면, 형식 $A$가 유형 $\rho$의 기본 수축 $\alpha_\rho$개로, 형식 $B$가 그와 같은 기본 수축 $\beta_\rho$개로 생겼다고 하자. 다음 경우에 $B$는 $A$보다 더 높다.
\begin{enumerate}
\renewcommand{\labelenumi}{\arabic{enumi})}
\item 부등식계 $\beta_\rho\geqq\alpha_\rho$, $\rho=1,2,\ldots,n-1$에서 적어도 하나의 $\rho$에 대해 엄격한 부등호가 성립하는 경우\srcfn{**)}{이 부등식계에서 일부에는 부호 $>$가, 일부에는 부호 $<$가 성립하면 두 형식은 서로 비교할 수 없다. 한 형식이 다른 형식에서 수축으로 생긴 경우에는 언제나 양의 기본 수축 값계가 대응하므로, 이 경우에는 $\beta_\rho-\alpha_\rho$가 양수이거나 0이어야 하기 때문이다.};
\item 모든 $\rho$에 대해 등호가 성립하지만 하나의 행렬곱으로 묶인 기호열의 수가 더 적은 경우. 이러한 정규화는 \S\ 8의 환원 때문에 필요함을 알 수 있다. 예를 들어 이에 따라 $\pair{S^{n-1}}{T_{n-1}}$은 $\pair{S^{n-1}}{[S^1T^\rho]_{n-\rho-1}p_\rho}$보다 더 높아진다.
\item 모든 $\rho$에 대해 등호가 성립하고 하나의 행렬곱으로 묶인 기호열의 수마저 같을 때에는, 그 자체로는 임의적이지만 한 번 정한 뒤에는 계속 유지하는 정규화에 의해 $B$를 $A$보다 더 높게 정한다. 하나의 값계 $\alpha$에 속하는 형식의 수는 유한하므로 이러한 정규화는 언제나 가능하다. 예를 들어 한 형식 $S$를 우선시하는 방식(기호열 $S$의 수가 더 많음, 열 $S$의 위쪽 무게지수 합이 더 큼 등)으로 이를 이룰 수 있다.
\end{enumerate}
형식열은 격자점들로 이루어진 $(n-1)$차원 다양체로 배열되어 있다고 생각하는 것이 편리하다. 이때 $n-1$개의 방향은 $n-1$개의 기본 수축에 대응한다. 그러면 각각의 형식에는 하나의 값계 $\alpha$, 곧 하나의 양의 정수 격자점이 유일하게 대응하고, 하나의 값계 $\alpha$에 속하는 여러 형식의 순서는 위의 정규화로 유일하게 정해진다. 형식열의 임의의 한 형식은 새로운 형식열의 출발형식이 되며, 이 새 형식열은 전체 형식열 안에 한 부분으로 포함된다. 더 정확히는 새 출발형식에서 시작하여 $n-1$개의 양의 방향으로 나아가 얻는 형식들의 부분으로 포함된다.

% inherited: noether_paper04_ko_translation_001_20260804\targets\Noether_P04_Korean_T09_U47_UNCHECKED.tex
% sha256: 9AAA9B2C6C361DE17EF42CC35DE661535482F1B4255EF916206BDE35FAAD4BBC
정리 VIII과 급수 전개의 일반 이론에 따라 이제 형식열에는 이항 및 삼항 영역의 환원 정리들이 그대로 성립한다(학위논문 \S\ 3 참조). 그 가운데 주정리를 간단히 유도하기에 앞서 다음 정의를 둔다.

“주어진 형식계에서 어떤 유한 개의 형식들이 하나의 모듈로 묶여 있다고 하자. 불변식을 인수로 갖는 형식들을 통해, 또는 ‘더 높은 형식들’을 통해 표현할 수 있는 형식을 \emph{환원 가능}이라 부른다. 여기서 더 높은 형식들이란 그 형식이 속하는 전체 형식열의 더 높은 형식들, 또는 그 모듈의 기호들을 더 높은 차수로 포함하는 형식들을 뜻한다. 환원 가능한 형식열을 \emph{환원자}라 부른다.” 그러면 다음이 성립한다.

\emph{정리 IX: 한 형식열의 출발형식이 어떤 환원자와의 수축으로 생겼기 때문에 환원 가능하다면, 이 출발형식에서 생기는 형식열 전체가 환원 가능하다.}

% inherited: noether_paper04_ko_translation_001_20260804\targets\Noether_P04_Korean_T09_U48_UNCHECKED.tex
% sha256: 8FE88B0477EEBCF7D038471D656F1AF0325060369FD7F5FDEFF8180AF8162F5F
증명을 위해, 형식 $f$와 또 다른 형식 $g$를 수축하여 생긴 형식 $h$는 (45.)에 의해 서로 동변인 여러 변수열 $p_\rho,p'_\rho$를 포함하는 형식 $f$로 볼 수 있음을 생각하자. 급수 전개의 일반 이론에 따르면 이러한 형식들은, 서로 동변인 변수열 $p_\rho,p'_\rho$ 두 개씩에 급수 전개를 차례로 적용함으로써 $f$의 기본 공변식들의 편극으로 나타낼 수 있다. 이때 (38.)에 따라 나타나는 기본 공변식들은 동변 수축으로 생성된 형식들이다. 그러나 급수 전개를 적용하는 순서는 여전히 임의적이므로, 특히 기본 수축들로부터 일반 수축을 생성하는 순서에 대응하는 순서를 택할 수 있다. 그러면 생기는 기본 공변식계는 형식열 $f$와 동일해지며, 결손이 더 높은 동변 수축으로 생긴 형식들은 기본 수축으로 생긴 형식들 속에 배열된다. 한편 급수 전개를 임의의 순서로 적용하여도 형식열 $f$의 편극들에 도달하며, 그 순서에는 두 번째 기본 공변식계가 대응한다고 하자. 실제로 형식열은 불변 구성식 전체를 소진하므로, 두 번째 계의 각 형식은 형식열의 편극들, 그중에서도 무게 감소가 같거나 더 큰 형식들의 편극들을 통해 선형으로 표현될 수 있다. 마지막 진술은 다음 사실에서 따른다. 곧 편극들은(\S\ 7 참조) 형식 $h$에 대응하는 차수들을 가진 형식열 $H$ (52.)와 형식열 $f$의 연립불변식으로 볼 수 있다. $H$를 자기 자신에 대해 수축할 때마다 하나의 무게 감소가 대응하며, 이 감소는 치환 (11.)을 시행한 뒤 기호열로, 따라서 $f$의 형식들로 옮겨간다\srcfn{*)}{이 고찰은 동치 개념의 두 번째 표현(\S\ 8 끝)의 바탕이기도 하다. 서로 다른 기본 공변식계에 관한 더 자세한 논의는 삼항 영역에 관한 Study, a. a. O., II. \S\ 7, 정리 7)과 8)에 있다. 우리의 정리 VII에 따라 같은 논의가 $n$개 변수에도 성립한다.}.

% inherited: noether_paper04_ko_translation_001_20260804\targets\Noether_P04_Korean_T09_U49_UNCHECKED.tex
% sha256: 3B4AAACCAAE9E1FB3E3FB99D99ADA033E36EC068AA6E34632DB6A14950284F48
형식열 $h$의 임의의 형식은 $h$를 자기 자신에 대해 수축하여 생긴다. 다시 말해 한편으로는 $g$와 $f$의 더 높은 수축을 통해, 다른 한편으로는 $f$ 또는 $g$를 자기 자신에 대해 수축하여 생긴다. 두 경우 모두 급수 전개는 출발형식 $h$에 관해 위에서 보인 것과 마찬가지로 형식열 $f$의 편극들로 이끈다. 특히 $f$가 환원자라면 환원자의 편극들에 따른 전개가 얻어지므로, 형식열 $h$는 환원 가능하다.

% inherited: noether_paper04_ko_translation_001_20260804\targets\Noether_P04_Korean_T09_U50_UNCHECKED.tex
% sha256: 218EF90531A8B0D9F2298E6BF340FC0D702848AD8BCF763E29D6A71C4DF0D036
완전 형식계의 환원에 이 정리를 적용하는 방법은 이항 및 삼항 영역에서와 마찬가지로 얻어진다.
